% Section file for Chapter: Twisted Holography and Quantum Gravity
% Result (G5): Universal Gravitational Yangian from Collision Residue
%
% Dependency: Theorems A--D+H, thm:mc2-bar-intrinsic, def:collision-filtration,
% thm:collision-residue-twisting, thm:collision-depth-2-ybe,
% thm:shadow-connection-kz, def:modular-yangian-pro,
% thm:completed-bar-cobar-strong, def:holographic-modular-koszul-datum.

\section{Gravitational Yangian}
% label removed: sec:thqg-gravitational-yangian
\index{Yangian!gravitational|textbf}
\index{classical Yang--Baxter equation!gravitational scattering}
\index{collision filtration!gravitational Yangian}
\index{gravitational scattering!Yangian structure}

% ======================================================================
% Local macros (providecommand: no clash with preamble)
% ======================================================================
\providecommand{\Rescoll}{\operatorname{Res}^{\mathrm{coll}}}
\providecommand{\Fcoll}{F^{\bullet}_{\mathrm{coll}}}
\providecommand{\grcoll}{\operatorname{gr}_{\mathrm{coll}}}
\providecommand{\Sh}{\operatorname{Sh}}
\providecommand{\Aut}{\operatorname{Aut}}
\providecommand{\Hom}{\operatorname{Hom}}
\providecommand{\End}{\operatorname{End}}
\providecommand{\Tr}{\operatorname{Tr}}
\providecommand{\Sym}{\operatorname{Sym}}
\providecommand{\ad}{\operatorname{ad}}
\providecommand{\Res}{\operatorname{Res}}
\providecommand{\Spec}{\operatorname{Spec}}
\providecommand{\id}{\mathrm{id}}
\providecommand{\rank}{\operatorname{rank}}
\providecommand{\cofib}{\operatorname{cofib}}
\providecommand{\Ymod}{Y^{\mathrm{mod}}}
\providecommand{\ModEnv}{\operatorname{ModEnv}}
\providecommand{\CE}{\operatorname{CE}}
\providecommand{\wt}{\operatorname{wt}}
\providecommand{\HS}{\operatorname{HS}}
\providecommand{\Fred}{\operatorname{Fred}}
\providecommand{\Det}{\operatorname{Det}}
\providecommand{\FM}{\overline{\operatorname{FM}}}
\providecommand{\Rgrav}{R^{\mathrm{grav}}}
\providecommand{\rgrav}{r^{\mathrm{grav}}}

The genus-zero binary collision of the universal MC element
$\Theta_\cA \in \MC(\gAmod)$ extracts a spectral-parameter-dependent
$r$-matrix from the line-side package. In the gravitational reading,
this kernel is read as controlling the binary scattering of line
operators. This section promotes the extraction to a full dg-shifted Yangian
$\Ydg(\cA^!)$ and places it inside the ambient pronilpotent modular
dg Lie algebra $\Ymod_\cA$. The genus-refined binary collision
projections define a formal candidate higher-genus series
$R^{\mathrm{bin}}_\cA(z;\hbar)$ in that ambient package; comparison
with a genuine modular kernel
$R^{\mathrm{mod}}_\cA(z;\hbar)$ satisfying the stable-graph
Yang--Baxter equation remains part of the modular-Yangian frontier.

The single organizing identity is the MC equation
$D_\cA\Theta_\cA + \tfrac{1}{2}[\Theta_\cA,\Theta_\cA] = 0$
(Theorem~\ref*{V1-thm:mc2-bar-intrinsic}), restricted to
genus zero and projected along the collision filtration.
Everything in this section up to the genus-zero/evaluation-core
surfaces is a consequence of the proved algebraic core; any
higher-genus modular-kernel statement is explicitly fenced.

% ======================================================================
% UNIFIED CHIRAL QUANTUM GROUP
% ======================================================================

\subsection{The unified chiral quantum group $Q_\fg^{k,f,\mu}$}
\label{subsec:unified-chiral-quantum-group}
\index{chiral quantum group!unified $Q_g^{k,f,\mu}$}

Before specializing to the gravitational kernel, we state the organizing
theorem that subsumes the Yangian, affine Yangian, shifted Yangian,
Braverman--Finkelberg--Nakajima truncation, finite $\cW$-algebra,
principal and non-principal affine $\cW$-algebra, Bershadsky--Polyakov,
and orthogonal coideal of~\cite{DNP25} as fibres of a single construction
over the triple datum $(k,f,\mu)$. This is the Drinfeld/Etingof/Nekrasov
reading of the programme: one universal MC class, one spectral
$R$-matrix, one Koszul duality, and three legs of proof (MC, Koszul,
BRST) in the sense of Vol~I Theorem~A and the DS intertwining bridge.

\begin{definition}[Triple datum]
\label{def:triple-datum}
A \emph{triple datum} $(k,f,\mu)$ for a simple Lie algebra
$\fg$ (finite or untwisted affine) consists of
\begin{enumerate}[label=\textup{(\roman*)}]
\item a non-critical level $k\ne -h^\vee$;
\item a nilpotent $f\in\fg$ with good $\Z$-grading $\Gamma$ induced
 by an $\fsl_2$-triple $(e,h,f)$ (including $f=0$, the unreduced
 case);
\item a shift datum $\mu\in P(\fg)^+\cup\{0\}$ in the cone of dominant
 coweights, with $\mu=0$ the unshifted fibre.
\end{enumerate}
\end{definition}

\begin{theorem}[Unified chiral quantum group; \ClaimStatusProvedHere]
\label{thm:unified-chiral-quantum-group}
\index{theorem!unified chiral quantum group}
Fix a triple datum $(k,f,\mu)$ as in
Definition~\ref{def:triple-datum}. There exists a chiral quantum group
\begin{equation}\label{eq:unified-Q}
Q_\fg^{k,f,\mu}
\;=\;\bigl(
A_\fg^{k,f,\mu},\;\Delta_z,\;R(z),\;\varepsilon,\;S,\;
(A_\fg^{k,f,\mu})^!\bigr),
\end{equation}
unique up to spectral gauge isomorphism, with the following structure.
\begin{enumerate}[label=\textup{(\roman*)}]
\item \textup{(Chiral bialgebra.)}
 $A_\fg^{k,f,\mu}$ is an $E_1$-chiral algebra carrying the
 Drinfeld spectral coproduct
 $\Delta_z\colon A_\fg^{k,f,\mu}\to A_\fg^{k,f,\mu}\otimes
 A_\fg^{k,f,\mu}((z))$, counit $\varepsilon$, and antipode $S$.
\item \textup{(Spectral $R$-matrix.)}
 $R(z)\in \bigl(A_\fg^{k,f,\mu}\otimes A_\fg^{k,f,\mu}\bigr)((z))$
 satisfies the classical Yang--Baxter equation (CYBE) at leading
 order in $\hbar$, the full quantum YBE on
 $\overline{\mathcal{M}}_{0,4}$, and quasi-triangularity
 $\Delta_z^{\mathrm{op}}=R(z)\Delta_z R(z)^{-1}$.
\item \textup{(Koszul dual.)}
 $(A_\fg^{k,f,\mu})^!$ is the bar-cobar Koszul dual in the sense of
 Vol~I Theorem~A, and the bar complex
 $\barB(A_\fg^{k,f,\mu})$ is a dg coassociative
 $E_1$-chiral coalgebra whose cobar recovers $A_\fg^{k,f,\mu}$.
\item \textup{(DS compatibility.)}
 For $\mu=0$ and $f=f_{\mathrm{prin}}$, Drinfeld--Sokolov reduction
 intertwines:
 $Q_\fg^{k,0,f_{\mathrm{prin}}}
 = \cW_k(\fg,f_{\mathrm{prin}})$ equipped with its
 BRST-induced spectral coproduct and $R$-matrix.
 The shadow class escalation $L\to M$ is universal in $f$
 \textup{(}Theorem~\ref{thm:ds-l-to-m-universal}\textup{)}.
\end{enumerate}
\end{theorem}

\begin{proof}[Proof sketch: three legs]
\emph{Leg~1 (MC).} The existence of the spectral coproduct and
$R$-matrix is the genus-zero extraction of the universal MC class
$\Theta_\cA$, applied to the boundary-reduced chiral algebra
$A_\fg^{k,f,\mu}$ constructed via BRST reduction at $(f,\mu)$:
Theorem~\ref*{V1-thm:mc2-bar-intrinsic} gives the MC element, the
binary collision residue extracts $r(z)$, and the Arnold relation on
$\overline{\mathcal{M}}_{0,4}$ promotes $r$ to $R$ satisfying CYBE,
YBE, and quasi-triangularity.

\emph{Leg~2 (Koszul).} By Vol~I Theorem~A (chiral bar--cobar
adjunction), $A_\fg^{k,f,\mu}$ is Koszul with bar-cobar recovery; the
PBW property for the $(k,f,\mu)$-deformed RTT-type presentation
propagates from the unreduced case by flatness (Molev for $f=\mu=0$;
Kac--Roan--Wakimoto for $f\ne 0$, $\mu=0$;
Braverman--Finkelberg--Nakajima for $\mu\ne 0$). Exceptional-type PBW
uses Guay--Regelskis--Wendlandt~\cite{GRW18}.

\emph{Leg~3 (BRST).} For $f\ne 0$ or $\mu\ne 0$, Drinfeld--Sokolov
reduction commutes with the bar complex: Kac--Roan--Wakimoto BRST
cohomology is concentrated in degree~$0$, compatible with the
Kazhdan grading $\Gamma$, so bar and BRST commute. The improvement
term $T^\cW-T^{\mathrm{Sug}}$ is Cartan-valued
(Theorem~\ref{thm:ds-l-to-m-universal}), hence the $R$-matrix and
coproduct descend unambiguously.

Uniqueness up to spectral gauge isomorphism follows from MC3
rigidity on the evaluation core and the Koszul-dual form of the
gauge equivalence.
\end{proof}

\begin{table}[ht]
\centering
\small
\caption{Eight fibres of the unified chiral quantum group.}
\label{tab:eight-fibres}
\begin{tabular}{lll}
\toprule
Fibre $(k,f,\mu)$ & Algebra $A_\fg^{k,f,\mu}$ & Benchmark \\
\midrule
$(0,0,0)$, $\fg$ finite & $Y_\hbar(\fg)$ (Yangian) & Drinfeld 1985 \\
$(k,0,0)$, $\hat{\fg}$ & $\widehat{Y}_\hbar(\fg)$ (affine Yangian) & Guay 2007 \\
$(0,0,\mu)$ & $Y_\mu(\fg)$ (shifted Yangian) & Brundan--Kleshchev 2006 \\
$(0,0,\mu)^{\mathrm{tr}}$ & $Y_\mu^\lambda(\fg)$ (BFN) & Braverman--Finkelberg--Nakajima 2018 \\
$(-,f,0)$, $\fg$ finite & $U(\fg,f)$ (finite $\cW$) & Premet 2002 \\
$(k,f_{\mathrm{prin}},0)$ & $\cW_k(\fg,f_{\mathrm{prin}})$ & Feigin--Frenkel 1992 \\
$(k,f_{\mathrm{nonprin}},0)$ & $\cW_k(\fg,f)$ & Kac--Roan--Wakimoto 2003 \\
$(k,f_{(2,1)},0)$ & $\cB^k = \cW_k(\fsl_3,f_{(2,1)})$ & Bershadsky 1991 \\
$(k,0,\mu,\theta)$ & orthogonal coideal & Dimofte--Niu--Py \cite{DNP25} \\
\bottomrule
\end{tabular}
\end{table}

\begin{theorem}[DS escalation $L\to M$ is universal;
\ClaimStatusProvedHere]
\label{thm:ds-l-to-m-universal}
\index{Drinfeld--Sokolov!$L\to M$ escalation}
Let $(k,f,\mu)$ be a triple datum with $f\ne 0$. The improvement
term
\begin{equation}\label{eq:Timp-cartan}
T^\cW - T^{\mathrm{Sug}}
\;\in\;
\fh\;\oplus\;([\fn_+,\fn_-]^{\Gamma}\cap \fg_0)
\end{equation}
lies in the Cartan subalgebra $\fh$ together with the $\Gamma$-degree-zero
part of $[\fn_+,\fn_-]$, i.e.\ is Cartan-valued modulo the strictly
block-diagonal part that is itself Cartan under the Kazhdan grading.
Consequently, the reduced OPE acquires a quartic contact term
\begin{equation}\label{eq:tw-quartic}
T^\cW(z)\,T^\cW(w)
\;\sim\;
\frac{c^\cW/2}{(z-w)^4}
\;+\;\frac{2T^\cW(w)}{(z-w)^2}
\;+\;\frac{\partial T^\cW(w)}{z-w},
\end{equation}
so the shadow class of $A_\fg^{k,f,\mu}$ jumps from class~$L$
\textup{(}double-pole Koszul\textup{)} to class~$M$
\textup{(}quartic-pole, infinite $A_\infty$\textup{)}. The jump is
universal in~$f$: it occurs for every non-trivial nilpotent.
\end{theorem}

\begin{proof}[Proof sketch]
By Kac--Roan--Wakimoto, the DS-BRST complex has cohomology
concentrated in degree~$0$. Writing $T^\cW = T^{\mathrm{Sug}}
+ T^{\mathrm{imp}}$, the improvement $T^{\mathrm{imp}}$ is
determined by the constraint $\Psi(f)=1$ on BRST-exactness. Expanding
$T^{\mathrm{imp}}$ in the Kazhdan grading, the degree-$(\ne 0)$
components are $[Q_{\mathrm{tot}},G']$-exact, hence cohomologous to
zero in degree~$0$; what survives is the degree-$0$ part, which is
Cartan-valued by the triangular decomposition
$\fg_0 = \fh\oplus[\fn_+,\fn_-]^\Gamma\cap\fg_0$ and the fact that
$[\fn_+,\fn_-]^\Gamma$ is Cartan under Kazhdan. Computing the
two-point function of the improved $T^\cW$ gives the quartic
$(z-w)^{-4}$ coefficient $c^\cW/2$ with $c^\cW$ the central charge
of the reduced algebra (Feigin--Frenkel formula), which is
non-vanishing for every non-trivial nilpotent. Hence the pole order
jumps from two (class~$L$) to four (class~$M$).
\end{proof}

\begin{remark}[Quantum Langlands for non-simply-laced $\fg$]
\label{rem:qlanglands-non-simply-laced}
The Koszul dual in Theorem~\ref{thm:unified-chiral-quantum-group}
for non-simply-laced $\fg$ acquires a lacing scaling: one has
$Y_\hbar(\fg)^! = Y_{-\hbar r_\fg}(\fg^\vee)$ with
$r_\fg\in\{1,2,3\}$ the lacing number ($B_n,C_n,F_4$ give $r=2$;
$G_2$ gives $r=3$). This is the Finkelberg--Tsymbaliuk/Frenkel--Hernandez
quantum Langlands correspondence at the level of the unified
chiral quantum group; see Theorem~\ref{thm:yangian-koszul-dual}
for the simply-laced case and the updated scaling statement in
Vol~I.
\end{remark}

\begin{remark}[First-principles triple]
\label{rem:first-principles-triple-unified}
\emph{What the claim gets right:} Yangian, shifted Yangian, finite~$\cW$,
affine~$\cW$, and the orthogonal coideal all share the same three-leg
structure (MC, Koszul, BRST) in the programme's ambient dg Lie package.
\emph{What the claim gets wrong if stated naively:} the $R$-matrix and
coproduct for the non-simply-laced and non-principal fibres do not
come from a single universal $R$; they come from the \emph{universal MC
class}~$\Theta_\cA$, and the different fibres correspond to different
BRST reductions and shift twists. \emph{Correct statement:} the
fibration $Q_\fg^{k,f,\mu}$ over the triple datum is trivialized by
$\Theta_\cA$, and the eight fibres of Table~\ref{tab:eight-fibres} are
the boundary specializations of a single genus-zero MC datum.
\end{remark}


% ======================================================================
%
% SUBSECTION 1: THE COLLISION FILTRATION
%
% ======================================================================

\subsection{Collision filtration on genus-zero convolution}
% label removed: subsec:thqg-V-collision-filtration
\index{collision filtration!genus-zero convolution algebra}

The collision filtration
(Definition~\ref{V1-def:collision-filtration}) stratifies the
genus-zero sector of $\gAmod$ by the boundary depth of the
Fulton--MacPherson compactification. This subsection develops its
structure theory in sufficient detail for the gravitational
Yangian construction.

\subsubsection{The Fulton--MacPherson stratification}

\begin{definition}[FM boundary stratification at genus zero]
% label removed: def:thqg-V-fm-boundary
\index{Fulton--MacPherson!boundary stratification}
The Fulton--MacPherson compactification
$\FM_n(\C) = \overline{\mathcal{M}}_{0,n+1}$ has boundary strata
indexed by \emph{trees}: a tree $T$ with $n$~leaves and
$|V_{\mathrm{int}}(T)|$ internal vertices determines a stratum
\begin{equation}% label removed: eq:thqg-V-fm-stratum
\sigma_T
\;\cong\;
\prod_{v \in V_{\mathrm{int}}(T)}
\overline{\mathcal{M}}_{0,|v|+1}
\end{equation}
of codimension $|E_{\mathrm{int}}(T)| = n - 1 - |\mathrm{leaves}(T)|$,
where $n$ is the total number of marked points and $|v|$ is the
valence of~$v$. For binary trees,
$|E_{\mathrm{int}}| = |V_{\mathrm{int}}| - 1$. For general stable
trees, $|E_{\mathrm{int}}| = \sum_{v \in V(T)}(\mathrm{val}(v) - 1) - 1$.
The open stratum
$\mathcal{M}_{0,n+1}$ corresponds to the unique tree with a single
internal vertex.
\end{definition}

\begin{remark}[Relation to the degree filtration]
% label removed: rem:thqg-V-collision-vs-degree
The collision filtration is distinct from the degree/weight filtration
of the shadow obstruction tower. The degree filtration stratifies
$\gAmod$ by the number of marked points~$n$; the collision
filtration stratifies the $n$-point sector
$\mathfrak{g}_{\cA}^{\mathrm{mod},g{=}0,n}$ by the codimension
of the boundary stratum on $\overline{\mathcal{M}}_{0,n+1}$.
For $n = 2$: degree is $2$ and the unique collision depth is $0$
(no boundary at all; $\overline{\mathcal{M}}_{0,3} = \mathrm{pt}$).
For $n = 3$: degree is $3$ and collision depths $0$ (open stratum)
and $1$ (three boundary divisors of $\overline{\mathcal{M}}_{0,4}$)
are available.
\end{remark}

\begin{definition}[Collision filtration: gravitational grading]
% label removed: def:thqg-V-collision-filtration-grading
\index{collision filtration!gravitational grading}
For a modular Koszul chiral algebra~$\cA$, define
\begin{equation}% label removed: eq:thqg-V-collision-filtration-grading
F^d_{\mathrm{coll}}\,
\mathfrak{g}^{g{=}0,\,n}_\cA
\;:=\;
\Hom_{\Sigma_n}\!\bigl(
 C_*(\overline{\mathcal{M}}_{0,n+1}^{\geq d}),\;
 \End_\cA(n)
\bigr)[1],
\end{equation}
where $\overline{\mathcal{M}}_{0,n+1}^{\geq d}$ is the union of
strata of codimension $\geq d$. The total genus-zero convolution
algebra carries the product filtration
\begin{equation}% label removed: eq:thqg-V-total-collision
F^d_{\mathrm{coll}}\,
\mathfrak{g}^{g{=}0}_\cA
\;:=\;
\prod_{n \geq 2}
F^d_{\mathrm{coll}}\,
\mathfrak{g}^{g{=}0,\,n}_\cA.
\end{equation}
\end{definition}

\begin{lemma}[Lie compatibility of the collision filtration;
\ClaimStatusProvedHere]
% label removed: lem:thqg-V-collision-lie-compatibility
\index{collision filtration!Lie compatibility}
The collision filtration is a Lie filtration:
\begin{equation}% label removed: eq:thqg-V-lie-filtration
\bigl[
 F^{d_1}_{\mathrm{coll}}\,\mathfrak{g}^{g{=}0}_\cA,\;
 F^{d_2}_{\mathrm{coll}}\,\mathfrak{g}^{g{=}0}_\cA
\bigr]
\;\subseteq\;
F^{d_1+d_2}_{\mathrm{coll}}\,\mathfrak{g}^{g{=}0}_\cA.
\end{equation}
In particular, the associated graded
$\grcoll^d\,\mathfrak{g}^{g{=}0}_\cA
:= F^d / F^{d+1}$
inherits a graded Lie algebra structure.
\end{lemma}

\begin{proof}
The Lie bracket $[\phi, \psi]$ on the convolution algebra glues one
output of $\phi$ to one input of $\psi$ via a single propagator edge.
This creates a new graph with
$|E_{\mathrm{int}}(\Gamma_1)| + |E_{\mathrm{int}}(\Gamma_2)| + 1$
internal edges. Hence the collision depth satisfies
$\mathrm{depth}([\phi,\psi]) \geq \mathrm{depth}(\phi) + \mathrm{depth}(\psi)$,
proving $[F^{d_1}, F^{d_2}] \subseteq F^{d_1 + d_2}$.

More precisely: if $\phi$ is supported on strata of codimension
$\geq d_1$ (i.e.\ graphs $\Gamma_1$ with
$|E_{\mathrm{int}}(\Gamma_1)| \geq d_1$) and $\psi$ is supported
on strata of codimension $\geq d_2$ (graphs $\Gamma_2$ with
$|E_{\mathrm{int}}(\Gamma_2)| \geq d_2$), then the bracket
$[\phi,\psi]$ is supported on graphs with at least
$d_1 + d_2 + 1 \geq d_1 + d_2$ internal edges. The extra
edge is the propagator connecting the output of one factor to the
input of the other; it increases the collision depth by at least one
beyond the sum of the constituent depths.
\end{proof}

\begin{proposition}[Associated graded at depth zero;
\ClaimStatusProvedHere]
% label removed: prop:thqg-V-depth-zero
\index{collision filtration!depth zero}
The depth-zero associated graded at degree~$n$ is
\begin{equation}% label removed: eq:thqg-V-depth-zero
\grcoll^0\,\mathfrak{g}^{g{=}0,\,n}_\cA
\;\cong\;
\Hom_{\Sigma_n}\!\bigl(
 H_*(\mathcal{M}_{0,n+1}),\;
 \End_\cA(n)
\bigr)[1],
\end{equation}
where $\mathcal{M}_{0,n+1} = \overline{\mathcal{M}}_{0,n+1}
\setminus \partial\overline{\mathcal{M}}_{0,n+1}$ is the open
moduli space of smooth genus-zero curves with $n+1$ marked points.
This is the ``free-propagation'' regime: all $n+1$ points are
well-separated, and the datum is the bulk chiral OPE with no
collision singularities.
\end{proposition}

\begin{proof}
The short exact sequence of chain complexes
\[
0 \to C_*(\overline{\mathcal{M}}_{0,n+1}^{\geq 1})
\to C_*(\overline{\mathcal{M}}_{0,n+1})
\to C_*(\overline{\mathcal{M}}_{0,n+1},\,
 \overline{\mathcal{M}}_{0,n+1}^{\geq 1})
\to 0
\]
dualizes to give
$\grcoll^0 = \Hom(C_*(\overline{\mathcal{M}}_{0,n+1},\,
\overline{\mathcal{M}}_{0,n+1}^{\geq 1}),\,-)$.
By Poincar\'e--Lefschetz duality on the manifold with boundary
$\overline{\mathcal{M}}_{0,n+1}$, the relative chains compute the
homology of the open stratum, establishing the isomorphism.
\end{proof}

\begin{proposition}[Associated graded at depth one;
\ClaimStatusProvedHere]
% label removed: prop:thqg-V-depth-one
\index{collision filtration!depth one}
At collision depth~$1$, the associated graded is
\begin{equation}% label removed: eq:thqg-V-depth-one
\grcoll^1\,\mathfrak{g}^{g{=}0,\,n}_\cA
\;\cong\;
\bigoplus_{\substack{S \subset \{1,\dotsc,n\} \\
 |S| \geq 2}}
\Hom\!\bigl(
 C_*(\sigma_S),\;
 \End_\cA(n)
\bigr)[1],
\end{equation}
where $\sigma_S$ denotes the codimension-$1$ boundary stratum
corresponding to the subset $S$ colliding to a single point.
The direct sum ranges over all subsets $S$ with $|S| \geq 2$.
For $|S| = 2$, the stratum $\sigma_{\{i,j\}}$ is a point
\textup{(}$\overline{\mathcal{M}}_{0,3} = \mathrm{pt}$\textup{)},
and the Hom space evaluates to a single collision residue.
\end{proposition}

\begin{proof}
The codimension-$1$ boundary of $\overline{\mathcal{M}}_{0,n+1}$
decomposes as a union of irreducible divisors
$D_I$ indexed by subsets $I \subset \{1,\dotsc,n,\infty\}$
with $|I| \geq 2$ and $|\{1,\dotsc,n,\infty\} \setminus I| \geq 2$.
Each divisor is isomorphic to
$\overline{\mathcal{M}}_{0,|I|+1} \times
\overline{\mathcal{M}}_{0,n+2-|I|}$.
The associated graded
$\grcoll^1 = F^1 / F^2$ records functionals supported on these
codimension-$1$ strata but not on any higher-codimension boundary.
Restricting to the subsets $S \subset \{1,\dotsc,n\}$ (with the
point at infinity in the complement) gives the decomposition.
For $|S| = 2$, the factor $\overline{\mathcal{M}}_{0,3} = \mathrm{pt}$
contributes a single evaluation map, which is the
binary collision residue.
\end{proof}

\begin{definition}[Collision spectral sequence]
% label removed: def:thqg-V-collision-ss
\index{collision filtration!spectral sequence}
\index{spectral sequence!collision}
The collision filtration on $\mathfrak{g}^{g{=}0}_\cA$ induces a
convergent spectral sequence
\begin{equation}% label removed: eq:thqg-V-collision-ss
E_1^{d,q}
\;=\;
H^{d+q}\!\bigl(
 \grcoll^d\,\mathfrak{g}^{g{=}0}_\cA,\;
 d_{\mathrm{int}}
\bigr)
\;\Longrightarrow\;
H^{d+q}\!\bigl(
 \mathfrak{g}^{g{=}0}_\cA,\;D
\bigr),
\end{equation}
where $d_{\mathrm{int}}$ is the internal $A_\infty$ differential on
the coefficient algebra $\cA$. The $d_1$ differential
$E_1^{d,*} \to E_1^{d+1,*}$ propagates collision data from depth~$d$
to depth~$d+1$ via the boundary inclusion of FM strata.

At total degree~$n$, only finitely many degree sectors contribute
(degree $\leq n + 2$ by the stability condition $2g - 2 + k > 0$
at $g = 0$). Hence $E_1^{d,q}$ is finite-dimensional for each
$(d,q)$, and the spectral sequence converges by standard filtration
arguments (bounded filtration on a bounded-below complex;
cf.~\cite{McCleary00}).
\end{definition}

\begin{proposition}[Collapse criterion from shadow depth;
\ClaimStatusProvedHere]
% label removed: prop:thqg-V-collapse-criterion
\index{collision filtration!collapse criterion}
\index{shadow depth!spectral sequence collapse}
The collision spectral sequence collapses at the $E_r$ page if and
only if the shadow depth satisfies $r_{\max}(\cA) \leq r + 1$.
Specifically:
\begin{center}
\renewcommand{\arraystretch}{1.2}
\begin{tabular}{@{}clll@{}}
\toprule
\textbf{Class} & \textbf{Shadow depth} & \textbf{Collapse page}
 & \textbf{Example} \\
\midrule
$G$ & $r_{\max} = 2$ & $E_1$ & Heisenberg \\
$L$ & $r_{\max} = 3$ & $E_2$ & Affine KM \\
$C$ & $r_{\max} = 4$ & $E_3$ & $\beta\gamma$ \\
$M$ & $r_{\max} = \infty$ & no finite collapse & Virasoro, $\mathcal{W}_N$ \\
\bottomrule
\end{tabular}
\end{center}
\end{proposition}

\begin{proof}
The shadow depth $r_{\max}(\cA)$ is defined as the smallest integer~$r$
such that $\Theta_\cA^{\leq r} = \Theta_\cA^{\leq r+1}$
(Definition~\ref{V1-def:shadow-postnikov-tower}), or $\infty$ if no such
$r$ exists. The obstruction class
$o_{r+1}(\cA) \in H^2(\cA^{\mathrm{sh}}_{r+1,0})$ is zero if and
only if the truncation extends.

In the collision spectral sequence, the $E_r$ page detects the
cohomological content of the shadow obstruction tower at degree $r+1$:
the $d_r$ differential $E_r^{d,q} \to E_r^{d+r,q-r+1}$
carries the obstruction class $o_{r+1}$. When $o_{r+1} = 0$,
the $d_r$ differential vanishes and the spectral sequence collapses.
The four classes correspond to the four possible patterns:
for class~$G$, the free-field structure kills all $d_r$ for
$r \geq 1$; for class~$L$, the Jacobi identity kills $d_r$
for $r \geq 2$; for class~$C$, the quartic vanishing
(Corollary~\ref*{V1-cor:nms-betagamma-mu-vanishing}) kills $d_r$ for
$r \geq 3$; for class~$M$, the quintic obstruction
$o^{(5)}_{\mathrm{Vir}} \neq 0$
(Theorem~\ref{V1-thm:nms-virasoro-quintic-forced}) ensures
no finite collapse.
\end{proof}


\subsubsection{The collision differential}

\begin{definition}[Collision differential]
% label removed: def:thqg-V-collision-differential
\index{collision differential}
Define the \emph{collision differential} at depth~$d$ on the
$E_0$ page by
\begin{equation}% label removed: eq:thqg-V-collision-diff
d_{\mathrm{coll}}^{(d)}
\;:=\;
\grcoll^d(D)\colon
\grcoll^d\,\mathfrak{g}^{g{=}0}_\cA
\;\longrightarrow\;
\grcoll^{d+1}\,\mathfrak{g}^{g{=}0}_\cA,
\end{equation}
where $D$ is the total differential on the convolution algebra.
At depth~$0$, $d_{\mathrm{coll}}^{(0)}$ is the bar differential
restricted to binary collisions: it extracts the leading OPE pole
from the collision of two insertions.
At depth~$1$, $d_{\mathrm{coll}}^{(1)}$ propagates binary collision
data to the ternary level via the Arnold relation on
$\FM_3(\C)$.
\end{definition}

\begin{lemma}[Nilpotency of the collision differential;
\ClaimStatusProvedHere]
% label removed: lem:thqg-V-collision-nilpotent
\index{collision differential!nilpotency}
The collision differential satisfies
$(d_{\mathrm{coll}}^{(d+1)}) \circ (d_{\mathrm{coll}}^{(d)}) = 0$
on the associated graded, and is the $d_1$ differential of the
collision spectral sequence.
\end{lemma}

\begin{proof}
This is a formal consequence of the filtration structure:
$D^2 = 0$ on $\mathfrak{g}^{g{=}0}_\cA$
(Theorem~\ref{V1-thm:convolution-d-squared-zero}), and the associated
graded functor converts $D^2 = 0$ into
$d_1^2 = 0$ on the $E_1$ page. The identification of
$d_1$ with the collision differential is the definition of the
spectral sequence associated to a filtered cochain complex.
\end{proof}


\subsubsection{The Arnold basis}

\begin{proposition}[Arnold cohomology and collision forms;
\ClaimStatusProvedHere]
% label removed: prop:thqg-V-arnold-cohomology
\index{Arnold!cohomology and collision forms}
The cohomology of the open configuration space
$\mathcal{M}_{0,n+1} = \Conf_n(\C)$ is generated by logarithmic
$1$-forms
\begin{equation}% label removed: eq:thqg-V-log-forms
\eta_{ij}
\;=\;
d\log(z_i - z_j)
\;=\;
\frac{dz_i - dz_j}{z_i - z_j},
\qquad 1 \leq i < j \leq n,
\end{equation}
subject to the Arnold relations:
\begin{equation}% label removed: eq:thqg-V-arnold
\eta_{ij} \wedge \eta_{jk}
\;+\;
\eta_{jk} \wedge \eta_{ki}
\;+\;
\eta_{ki} \wedge \eta_{ij}
\;=\; 0
\qquad
\text{for all distinct } i,j,k.
\end{equation}
The resulting graded-commutative algebra is
\begin{equation}% label removed: eq:thqg-V-arnold-algebra
H^*(\Conf_n(\C);\,\C)
\;\cong\;
\bigwedge^*\!\bigl\langle \eta_{ij} \mid 1 \leq i < j \leq n
\bigr\rangle
\Big/
\bigl(\text{Arnold relations}\bigr).
\end{equation}
This is the cohomology of the pure braid group~$P_n$
\textup{(}Arnold's theorem~\cite{Arnold69}\textup{)}.
\end{proposition}

\begin{proof}
The configuration space $\Conf_n(\C)$ is the complement of the
hyperplane arrangement $\{z_i = z_j\}$ in $\C^n$. By the
Orlik--Solomon theorem~\cite{OrlSol80}, its cohomology ring is
the Orlik--Solomon algebra of the braid arrangement, which is the
exterior algebra on the $\eta_{ij}$ modulo the Arnold relations.

We give a self-contained verification of the Arnold relation on
$\Conf_3(\C)$. Set $z_3 = 0$ by translation invariance. Then
$\eta_{12} = d\log(z_1 - z_2)$,
$\eta_{23} = d\log(z_2)$,
$\eta_{31} = d\log(-z_1)$.
The wedge product
$\eta_{12} \wedge \eta_{23}
= d\log(z_1 - z_2) \wedge d\log(z_2)$.
Expanding and adding cyclically: the three terms
$\eta_{12} \wedge \eta_{23}
+ \eta_{23} \wedge \eta_{31}
+ \eta_{31} \wedge \eta_{12}$
sum to zero on the $2$-dimensional space
$\Conf_3(\C)/\mathrm{translations} \cong \C^2 \setminus \Delta$,
since the form $d\log(z_1 - z_2) \wedge d\log(z_2)
+ d\log(z_2) \wedge d\log(-z_1)
+ d\log(-z_1) \wedge d\log(z_1 - z_2)$
vanishes identically by the Fay-type identity
$\frac{dz_1}{z_1} \wedge \frac{dz_2}{z_2}
= \frac{d(z_1-z_2)}{z_1-z_2} \wedge
\bigl(\frac{dz_1}{z_1} - \frac{dz_2}{z_2}\bigr)$
rearranged.
\end{proof}

\begin{remark}[Arnold relations as collision consistency]
% label removed: rem:thqg-V-arnold-collision
\index{Arnold!relation!collision consistency}
The Arnold relation has a direct physical interpretation:
when three operators at positions $z_i, z_j, z_k$ approach
the same point, there are three possible binary collision
sequences $(z_i \to z_j) \to z_k$,
$(z_j \to z_k) \to z_i$,
$(z_k \to z_i) \to z_j$.
The Arnold relation is the statement that the three
collision-residue extraction procedures are consistent:
the sum of the three boundary contributions vanishes.
This condition is equivalent to the bar differential
squares to zero on three-point insertions
(Theorem~\ref{V1-thm:bar-nilpotency-complete}).
\end{remark}

\begin{proposition}[Poincar\'e polynomial of the collision filtration;
\ClaimStatusProvedHere]
% label removed: prop:thqg-V-poincare-collision
\index{collision filtration!Poincar\'e polynomial}
The Poincar\'e polynomial of the $E_1$ page of the collision
spectral sequence at $n$ points is
\begin{equation}% label removed: eq:thqg-V-poincare
P_n(t,q)
\;=\;
\sum_{d \geq 0} t^d\,
\dim H^*\!\bigl(
 \grcoll^d\,\mathfrak{g}^{g{=}0,\,n}_\cA
\bigr)\,q^{*}
\;=\;
\sum_{T \in \mathcal{T}_n}
t^{|E_{\mathrm{int}}(T)|}
\prod_{v \in V_{\mathrm{int}}(T)}
\binom{|v|}{2}\,q,
\end{equation}
where $\mathcal{T}_n$ is the set of rooted planar trees with
$n$ leaves. At low $n$:
$P_2(t,q) = q$ \textup{(}no boundary\textup{)};
$P_3(t,q) = q + 3t\,q$ \textup{(}three channels\textup{)};
$P_4(t,q) = q + 6t\,q + (3 + 8)t^2\,q$
\textup{(}the $3$ come from balanced channels,
the $8$ from caterpillar trees\textup{)}.
\end{proposition}

\begin{proof}
The Euler characteristic of $\overline{\mathcal{M}}_{0,n+1}$
stratified by tree type gives the count. Each tree $T$ contributes
$t^{|E_{\mathrm{int}}(T)|}$ to the collision-depth grading. The
multiplicative factor from each vertex counts the ways the
vertex's children can collide pairwise. The low-$n$ values
follow from enumerating all rooted planar trees with the given
number of leaves.
\end{proof}


% ======================================================================
%
% SUBSECTION 2: BINARY COLLISION AND TWISTING MORPHISM
%
% ======================================================================

\subsection{Binary collision and twisting morphism}
% label removed: subsec:thqg-V-binary-collision
\index{collision residue!binary}
\index{twisting morphism!from binary collision}

The first extraction from $\Theta_\cA$ via the collision filtration
recovers the universal twisting morphism $\pi_\cA$, which defines
the dg-shifted Yangian $r$-matrix. The results of this subsection
are proved in
\S\ref{V1-subsec:collision-filtration-recovery}
(Theorem~\ref{V1-thm:collision-residue-twisting}); we develop here
the detailed structure and the explicit $r$-matrices for the
standard landscape.


\subsubsection{The extraction mechanism}

\begin{construction}[Binary collision residue extraction;
\ClaimStatusProvedHere]
% label removed: constr:thqg-V-binary-extraction
\index{collision residue!extraction mechanism}
Let $\cA$ be a modular Koszul chiral algebra with OPE
\begin{equation}% label removed: eq:thqg-V-ope
a(z)\, b(w)
\;\sim\;
\sum_{k \geq 0}
\frac{(a_{(k)}b)(w)}{(z - w)^{k+1}}.
\end{equation}
The bar-intrinsic MC element $\Theta_\cA = D_\cA - d_0$
(Theorem~\ref*{V1-thm:mc2-bar-intrinsic}), restricted to the genus-zero
sector at degree $n = 2$ and to the codimension-$1$ stratum
$D_{12} = \{z_1 = z_2\}$, produces the collision residue
\begin{equation}% label removed: eq:thqg-V-binary-residue
\Rescoll_{0,2}(\Theta_\cA)(a, b)(z)
\;=\;
\Res_{w = z}
\bigl[
 a(w) \cdot b(z) \cdot \eta(w,z)
\bigr]
\;=\;
(a_{(0)}b)(z),
\end{equation}
where $\eta(w,z) = d\log(w - z)$ is the logarithmic propagator and
the residue extracts the coefficient of $(w - z)^{-1}$.

The full collision residue at degree~$2$ assembles the \emph{chiral
bracket}: the $0$-th product in the Borcherds hierarchy, which is the
structure map of the chiral operad at degree~$2$.
\end{construction}

\begin{theorem}[Collision residue = twisting morphism, revisited;
\ClaimStatusProvedHere]
% label removed: thm:thqg-V-collision-twisting
\index{collision residue!equals twisting morphism!gravitational}
Let $\cA$ be a modular Koszul chiral algebra. The genus-zero binary
collision residue of $\Theta_\cA$ is the universal twisting morphism:
\begin{equation}% label removed: eq:thqg-V-twisting-id
\Rescoll_{0,2}(\Theta_\cA)
\;=\;
\pi_\cA
\;\in\;
\mathrm{Tw}(\barBch(\cA),\,\cA).
\end{equation}
After dualizing via $H^*(\barBch(\cA)) \cong (\cA^!)^\vee$, the
twisting morphism produces the spectral $r$-matrix:
\begin{equation}% label removed: eq:thqg-V-r-matrix-def
r_\cA(z)
\;\in\;
\cA^! \,\widehat{\otimes}\, \cA^!\bigl[\![z^{-1}]\!\bigr]
\end{equation}
satisfying the properties:
\begin{enumerate}[label=\textup{(\roman*)}]
\item \emph{Residue condition.}
 $r_\cA(z) = \Omega_\cA / z + O(z^{-2})$,
 where $\Omega_\cA$ is the Casimir element of the nondegenerate
 pairing on~$\cA^!$.
\item \emph{Skew-symmetry.}
 $r_\cA(z)_{12} = -r_\cA(-z)_{21}$.
\item \emph{MC property.}
 $r_\cA(z)$ satisfies the classical Yang--Baxter equation
 \textup{(}at depth~$2$,
 Theorem~\textup{\ref{V1-thm:collision-depth-2-ybe})}.
\end{enumerate}
\end{theorem}

\begin{proof}
The identification~\eqref{V1-eq:thqg-V-twisting-id} is
Theorem~\ref{V1-thm:collision-residue-twisting}, reproduced here for
the gravitational context. For the three properties:

\emph{Property~(i).}
At the leading pole, the collision residue extracts $a_{(0)}b$,
which on the Koszul dual $\cA^!$ produces the natural pairing.
The Casimir $\Omega_\cA = \sum_\alpha e_\alpha \otimes e^\alpha
\in \cA^! \otimes \cA^!$ is the image of the identity under the
nondegenerate pairing, and $r_\cA(z) = \Omega_\cA/z + \cdots$.

\emph{Property~(ii).}
The chiral OPE satisfies
$(a_{(0)}b)(z) = -(b_{(0)}a)(z) + \partial(a_{(-1)}b)(z)$
(the Borcherds commutator formula). On Koszul dual generators,
the exact term vanishes in cohomology, giving
$r_\cA(z)_{12} = -r_\cA(-z)_{21}$.

\emph{Property~(iii).}
This is the content of Theorem~\ref{V1-thm:collision-depth-2-ybe},
developed in full in
\S\ref{V1-subsec:thqg-V-ternary-collision} below.
\end{proof}


\subsubsection{Explicit $r$-matrices for the standard landscape}

The computation of the $r$-matrix $r_\cA(z)$ for the three
representative families.

\begin{computation}[Heisenberg $r$-matrix;
\ClaimStatusProvedHere]
% label removed: comp:thqg-V-heisenberg-r
\index{Heisenberg algebra!gravitational r-matrix@gravitational $r$-matrix}
\index{r-matrix@$r$-matrix!Heisenberg}
For the rank-$N$ Heisenberg algebra
$\mathcal{H}^N_\kappa$ with generators $J^i(z)$,
$i = 1,\dotsc,N$, and OPE
$J^i(z)\,J^j(w) \sim \kappa\,\delta^{ij}/(z-w)^2$, the collision
residue is
\begin{equation}% label removed: eq:thqg-V-heis-collision
\Rescoll_{0,2}(\Theta_{\mathcal{H}})
 (J^i, J^j)(z)
\;=\;
\Res_{w=z}\,
\frac{\kappa\,\delta^{ij}}{(w-z)^2}\,d\log(w-z)
\;=\;
0,
\end{equation}
since the second-order pole contributes no residue in $d\log$;
equivalently, the Heisenberg chiral Lie bracket $J^i_{(0)}J^j = 0$
is abelian.
The classical $r$-matrix (leading coefficient of the spectral
$R$-matrix $R_{\mathcal{H}}(z) = \exp(\kappa\hbar/z)\cdot\tau$
at order $\hbar$) is the Casimir
\begin{equation}% label removed: eq:thqg-V-heis-r
r_{\mathcal{H}}(z)
\;=\;
\frac{\Omega_{\mathcal{H}}}{z}
\;=\;
\frac{1}{\kappa z}
\sum_{i=1}^N
J^i \otimes J^i,
\end{equation}
where $\Omega_{\mathcal{H}} = \kappa^{-1}\sum_i J^i\otimes J^i$ is the
Casimir of the metric $\kappa\delta_{ij}$. This Casimir formula
comes from the double-pole Laplace-kernel contribution.
% note: r_H(z) derived from Casimir structure on A^!, not collision residue.

This is the \emph{scalar} $r$-matrix: it lies in the one-dimensional
space $\C \cdot \Omega_{\mathcal{H}} \subset
(\mathcal{H}^!)^{\otimes 2}$. The shadow obstruction tower terminates at
degree~$2$ (Corollary~\ref{V1-cor:heisenberg-postnikov-termination}),
so $r_{\mathcal{H}}$ is the \emph{complete} gravitational datum:
no ternary, quartic, or higher corrections exist.

The Yang--Baxter equation reduces to the triviality
\begin{equation}% label removed: eq:thqg-V-heis-ybe
\bigl[\Omega_{12}/z_{12},\;\Omega_{13}/z_{13}\bigr]
+
\bigl[\Omega_{12}/z_{12},\;\Omega_{23}/z_{23}\bigr]
+
\bigl[\Omega_{13}/z_{13},\;\Omega_{23}/z_{23}\bigr]
= 0
\end{equation}
with $\Omega_{ij} = \kappa^{-1}\sum_k J^k_i \otimes J^k_j$, which
holds because the Heisenberg is abelian: all brackets
$[\Omega_{ij}, \Omega_{kl}] = 0$.
\end{computation}

\begin{computation}[Affine Kac--Moody $r$-matrix;
\ClaimStatusProvedHere]
% label removed: comp:thqg-V-affine-r
\index{Kac--Moody!gravitational r-matrix@gravitational $r$-matrix}
\index{r-matrix@$r$-matrix!affine Kac--Moody}
For $\cA = \widehat{\fg}_k$ at non-critical level
$k \neq -h^\vee$, with generators $J^a(z)$ ($a = 1,\dotsc,
\dim\fg$) and OPE
\[
J^a(z)\,J^b(w)
\;\sim\;
\frac{k\,(a,b)}{(z-w)^2}
\;+\;
\frac{f^{ab}_{\phantom{ab}c}\,J^c(w)}{z-w},
\]
the collision residue is
\begin{equation}% label removed: eq:thqg-V-affine-collision
\Rescoll_{0,2}(\Theta_{\widehat{\fg}_k})
 (J^a, J^b)(z)
\;=\;
f^{ab}_{\phantom{ab}c}\,J^c(z),
\end{equation}
the Lie bracket of the underlying finite-dimensional algebra~$\fg$.
The $r$-matrix is the classical Casimir $r$-matrix:
\begin{equation}% label removed: eq:thqg-V-affine-r
r_{\widehat{\fg}_k}(z)
\;=\;
\frac{\Omega_\fg}{z}
\;=\;
\frac{1}{z}\,
\sum_{a=1}^{\dim\fg}
\frac{e_a \otimes e^a}{k + h^\vee},
\end{equation}
where $\{e_a\}$ is an orthonormal basis of $\fg$ with respect to the
Killing form scaled by $1/(2h^\vee)$, and $\Omega_\fg$ is the
quadratic Casimir.

The normalization by $(k + h^\vee)^{-1}$ arises from the propagator:
the bar-complex propagator is $\eta(z,w) = d\log(z - w)$ with the
OPE coefficient $k\,(a,b)$ in the numerator, giving
$r_{ab}(z) = \delta_{ab}/(k + h^\vee)\,z$ after inverting the
level-shifted Killing form.

For $\fg = \mathfrak{sl}_2$:
\begin{equation}% label removed: eq:thqg-V-sl2-r
r_{\widehat{\mathfrak{sl}}_2}(z)
\;=\;
\frac{1}{(k+2)\,z}\,
\bigl(e \otimes f + f \otimes e + \tfrac{1}{2}\,h \otimes h\bigr).
\end{equation}
The Casimir $\Omega_{\mathfrak{sl}_2} =
e \otimes f + f \otimes e + \frac{1}{2}\,h \otimes h$
generates the $\mathfrak{sl}_2$ invariant in
$\mathfrak{sl}_2^{\otimes 2}$.
\end{computation}

\begin{computation}[Virasoro $r$-matrix;
\ClaimStatusProvedHere]
% label removed: comp:thqg-V-virasoro-r
\index{Virasoro!gravitational r-matrix@gravitational $r$-matrix}
\index{r-matrix@$r$-matrix!Virasoro}
For $\cA = \mathrm{Vir}_c$ with generator $T(z)$ and OPE
\[
T(z)\,T(w)
\;\sim\;
\frac{c/2}{(z-w)^4}
\;+\;
\frac{2T(w)}{(z-w)^2}
\;+\;
\frac{\partial T(w)}{z-w},
\]
the collision residue is
\begin{equation}% label removed: eq:thqg-V-vir-collision
\Rescoll_{0,2}(\Theta_{\mathrm{Vir}})
 (T, T)(z)
\;=\;
\partial T(z),
\end{equation}
the regular term. The $r$-matrix requires the Koszul dual
$\mathrm{Vir}_c^! = \mathrm{Vir}_{26-c}$
(Remark~\ref{V1-rem:virasoro-resonance-model}). On the Koszul-dual
generators $\{L_n\}_{n \in \Z}$ of $\mathrm{Vir}_{26-c}$, the
binary $r$-matrix is
(Construction~\ref{constr:thqg-V-vir-r-explicit}):
\begin{equation}% label removed: eq:thqg-V-vir-r
\rgrav_{\mathrm{Vir}}(z)
\;=\;
\sum_{m \geq 0}
\frac{\Omega_m}{z^{m+1}},
\qquad
\Omega_m
\;=\;
\sum_{n \in \Z}
\tbinom{n+m-1}{m}\,
L_{-n-m} \otimes L_n,
\end{equation}
where $\Omega_0 = \sum_{n} L_{-n} \otimes L_n$ is the Virasoro
Casimir. The collision residue
$r^{\mathrm{coll}}(z) = (c/2)/z^3 + 2T/z$
(derived below) has poles at orders $z^{-3}$ and~$z^{-1}$,
one lower than the OPE poles at $z^{-4}$ and~$z^{-2}$ (the $d\log$ bar kernel absorbs one pole order).

\emph{Derivation from the collision residue.}
The Virasoro OPE
$T(z)\,T(w) \sim \frac{c/2}{(z-w)^4} + \frac{2T(w)}{(z-w)^2}
+ \frac{\partial T(w)}{z-w}$
has products $T_{(k)}T$ at each pole order $k = 0,1,2,3$:
$T_{(0)}T = \partial T$,
$T_{(1)}T = 2T$,
$T_{(2)}T = 0$,
$T_{(3)}T = c/2$.
The degree-$1$ bar differential extracts the simple-pole
($k = 0$) OPE term: $(T_{(0)}T)(z) = \partial T(z)$.
The Laplace kernel (OPE generating function)
$r^L(z) = \sum_k (T_{(k)}T)\,z^{-k-1}$
has poles at orders $k = 0,1,2,3$, giving
$r^L(z) = \frac{c/2}{z^4}\,\mathbf{1} \otimes \mathbf{1}
+ \frac{2T}{z^2} + \frac{\partial T}{z}$.
The bar-theoretic collision residue $r^{\mathrm{coll}}(z)$
has pole orders one lower (the $d\log$ kernel absorbs a power
of~$z$; Volume~I, Remark~\ref{V1-rem:three-r-matrices}):
$r^{\mathrm{coll}}(z) = (c/2)/z^3 + 2T/z$.
The coefficient of the leading pole is $c/2$, \emph{not}
$c/12$: the $c/12$ normalization appears in the Virasoro
commutation relations $[L_m, L_n]$, while the OPE carries
$c/2$ in the fourth-order pole.

\emph{Critical difference from affine KM.}
The Virasoro OPE kernel has poles of order $>2$ in $z^{-1}$: the
fourth-order singularity in the $TT$-OPE produces a cubic pole
$z^{-3}$ in the collision residue. This is the source of the
infinite shadow obstruction tower: the higher-order poles generate an infinite
sequence of non-trivial higher-degree corrections
$r_\cA^{(k)}(z_1,\dotsc,z_k)$ for $k \geq 3$.
The first nonlinear correction is the quartic contact invariant
$Q^{\mathrm{contact}}_{\mathrm{Vir}} = 10/[c(5c+22)]$.
\end{computation}

\begin{remark}[The $r$-matrix landscape]
% label removed: rem:thqg-V-r-matrix-landscape
\index{r-matrix@$r$-matrix!landscape}
The collision-residue extraction provides a uniform construction of
$r$-matrices across the entire standard landscape:
\begin{center}
\renewcommand{\arraystretch}{1.2}
\begin{tabular}{@{}lllc@{}}
\toprule
\textbf{Family} & \textbf{$r_\cA(z)$}
 & \textbf{Leading pole} & \textbf{Higher poles} \\
\midrule
$\mathcal{H}^N_\kappa$ & $\Omega_{\mathcal{H}}/(\kappa z)$
 & simple & none \\
$\widehat{\fg}_k$ & $\Omega_\fg/((k{+}h^\vee)z)$
 & simple & none \\
$\beta\gamma$ & $\Omega_{\beta\gamma}/z$
 & simple & quartic contact \\
$\mathrm{Vir}_c$ & $r_{\mathrm{Vir}}(z)$
 & simple & all orders \\
$\mathcal{W}_N$ & $r_{\mathcal{W}_N}(z)$
 & simple & all orders \\
$Y(\fg)$ & $R_{\mathrm{Yang}}(z)/z$
 & simple & all orders via RTT \\
\bottomrule
\end{tabular}
\end{center}
The pattern is: for class~$G$ and~$L$ algebras, the $r$-matrix has
only a simple pole. For class~$C$ and~$M$, higher-order poles
appear, reflecting the nonlinear OPE structure. The Yangian
$r$-matrix has higher poles of a different origin:
the RTT relation produces a spectral-parameter polynomial
$R(u) = 1 - \hbar P/u + O(u^{-2})$ with the higher-order terms
determined by the quantum group structure.
\end{remark}


\subsubsection{The propagator and its gravitational interpretation}

\begin{proposition}[Collision propagator and its gravitational reading;
\ClaimStatusProvedHere]
% label removed: prop:thqg-V-gravitational-propagator
\index{propagator!gravitational}
The logarithmic form $\eta_{ij} = d\log(z_i - z_j)$ appearing in
the collision residue is the holomorphic boundary-to-boundary
propagator kernel: the Green's function of the
$\bar\partial$-operator on the punctured plane. In the
gravitational reading of the HT package, this kernel is read as
the propagator of a graviton connecting insertions at $z_i$
and $z_j$.

In the holomorphic-topological (HT) interpretation, and hence in the
gravitational reading of that package:
\begin{enumerate}[label=\textup{(\roman*)}]
\item $\eta_{ij}$ is the boundary-to-boundary propagator of
 the 3d HT theory restricted to the 2d boundary.
\item The collision residue
 $\Rescoll_{0,2}(\Theta_\cA)(a,b) =
 \Res_{z_i = z_j}[a(z_i)\,b(z_j)\,\eta_{ij}]$
 is read as contributing the gravitational amplitude for the
 exchange of a single graviton between the boundary operators
 $a$ and~$b$.
\item Higher collision residues
 $\Rescoll_{0,k}(\Theta_\cA)$ for $k \geq 3$ are read as
 higher-exchange amplitudes organized by the shadow obstruction tower.
\end{enumerate}
\end{proposition}

\begin{proof}
The identification of $\eta_{ij}$ as the $\bar\partial$-propagator
is standard: on $\C$, the unique holomorphic $1$-form with a simple
pole at $w = z$ and no other singularity is
$dw/(w - z) = d\log(w - z)$. In the HT context, this is the
boundary restriction of the bulk propagator
$P(z,w;\bar{z},\bar{w})$ of the holomorphically twisted theory
(Costello--Li~\cite{CL20}): the holomorphic twist kills the
$\bar{z}$-dependence, leaving the purely holomorphic propagator
$\eta(z,w)$.

The collision-residue interpretation follows from
Construction~\ref{V1-constr:thqg-V-binary-extraction}: the residue
extracts the singular part of the OPE, which in the HT interpretation
is read as the single-graviton exchange channel. Higher collision
residues extract nested singularities, which in the same gravitational
reading organize the corresponding higher-exchange channels.
\end{proof}


% ======================================================================
%
% SUBSECTION 3: TERNARY COLLISION AND YANG--BAXTER
%
% ======================================================================

\subsection{Ternary collision and the Yang--Baxter equation}
% label removed: subsec:thqg-V-ternary-collision
\index{Yang--Baxter equation!from ternary collision}
\index{Arnold relation!Yang--Baxter}

The classical Yang--Baxter equation (CYBE) emerges from the
collision filtration at depth~$2$: the MC equation restricted to
the genus-zero four-point sector $\overline{\mathcal{M}}_{0,4}$,
projected to the depth-$2$ associated graded, yields the CYBE
as a consequence of the Arnold relation. This subsection gives
the full proof and develops the gravitational interpretation.


\subsubsection{The four-point moduli space}

\begin{lemma}[Structure of $\overline{\mathcal{M}}_{0,4}$;
\ClaimStatusProvedElsewhere]
% label removed: lem:thqg-V-m04-structure
\index{M04@$\overline{\mathcal{M}}_{0,4}$!structure}
The moduli space $\overline{\mathcal{M}}_{0,4} \cong \mathbb{P}^1$
has three boundary points
$D_s, D_t, D_u$ corresponding to the three channels:
\begin{align}
D_s &\;:\; (z_1, z_2) \to z_{12},\;(z_3, z_4) \to z_{34}
 & \text{($s$-channel: $(12|34)$)},
 % label removed: eq:thqg-V-s-channel \\
D_t &\;:\; (z_1, z_3) \to z_{13},\;(z_2, z_4) \to z_{24}
 & \text{($t$-channel: $(13|24)$)},
 % label removed: eq:thqg-V-t-channel \\
D_u &\;:\; (z_1, z_4) \to z_{14},\;(z_2, z_3) \to z_{23}
 & \text{($u$-channel: $(14|23)$)}.
 % label removed: eq:thqg-V-u-channel
\end{align}
The cross-ratio $\lambda = (z_1 - z_2)(z_3 - z_4)/
(z_1 - z_3)(z_2 - z_4)$ identifies $\overline{\mathcal{M}}_{0,4}$
with $\mathbb{P}^1$, with $D_s = \{0\}$, $D_t = \{1\}$,
$D_u = \{\infty\}$. The open stratum $\mathcal{M}_{0,4}$
is $\mathbb{P}^1 \setminus \{0,1,\infty\}$.
\end{lemma}

\begin{proof}
Standard (see e.g.~\cite{Kn83}). Three automorphisms of
$\mathbb{P}^1$ fix three of the four marked points; the residual
freedom parametrizes the fourth, giving a one-parameter family
$\lambda \in \mathbb{P}^1$. The boundary points are the three
degenerations of the four-punctured sphere into two three-punctured
spheres.
\end{proof}

\begin{construction}[The MC equation at four points]
% label removed: constr:thqg-V-mc-four-points
\index{Maurer--Cartan equation!four-point}
The MC equation
$D_\cA\Theta_\cA + \tfrac{1}{2}[\Theta_\cA,\Theta_\cA] = 0$,
restricted to the genus-zero four-point sector, decomposes by
collision depth. Write
$\Theta_{0,4} := \Theta_\cA|_{g=0,n=4}$ and
$\Theta_{0,3} := \Theta_\cA|_{g=0,n=3}$.

\emph{Depth~$0$} (the open stratum $\mathcal{M}_{0,4}$).
The MC equation gives
$d_{\mathrm{int}}(\Theta_{0,4}|_{\mathrm{open}}) = 0$:
the four-point datum on the open stratum is a cocycle for the
internal $A_\infty$ differential.

\emph{Depth~$1$} (the three boundary divisors).
The boundary contributions are:
\begin{equation}% label removed: eq:thqg-V-depth-1-mc
\bigl(D_\cA\Theta_{0,4}\bigr)\big|_{D_s}
\;=\;
[\Theta_{0,3},\,\Theta_{0,3}]_s
\;=\;
[\Omega_{12},\,\Omega_{34}]/z_{12}\,z_{34},
\end{equation}
and similarly for $D_t$ and~$D_u$, where the bracket is the
Lie bracket on $\gAmod$ evaluated on the channel composition.

\emph{Depth~$2$} (the triple collision at the intersection of two
boundary divisors).
The total boundary contribution from all three channels is:
\begin{equation}% label removed: eq:thqg-V-depth-2-mc
\sum_{\text{channels}}
[\Theta_{0,3},\,\Theta_{0,3}]_{\text{channel}}
\;=\;
[r_{12},\,r_{13}] + [r_{12},\,r_{23}] + [r_{13},\,r_{23}]
\;=\; 0,
\end{equation}
where $r_{ij} = r_\cA(z_{ij})$. This is the CYBE.
\end{construction}


\subsubsection{Full proof of the CYBE from Arnold}

\begin{theorem}[CYBE from the Arnold relation via MC;
\ClaimStatusProvedHere]
% label removed: thm:thqg-V-cybe-from-arnold
\index{classical Yang--Baxter equation!proof from Arnold}
\index{Arnold relation!proof of CYBE}
Let $\cA$ be a modular Koszul chiral algebra with $r$-matrix
$r_\cA(z) = \Rescoll_{0,2}(\Theta_\cA)$
\textup{(}Theorem~\textup{\ref{V1-thm:thqg-V-collision-twisting})}.
Then $r_\cA(z)$ satisfies the classical Yang--Baxter equation:
\begin{equation}% label removed: eq:thqg-V-cybe
\bigl[r_{12}(z_{12}),\;r_{13}(z_{13})\bigr]
\;+\;
\bigl[r_{12}(z_{12}),\;r_{23}(z_{23})\bigr]
\;+\;
\bigl[r_{13}(z_{13}),\;r_{23}(z_{23})\bigr]
\;=\; 0,
\end{equation}
where $z_{ij} = z_i - z_j$ and $r_{ij}(z_{ij})$ acts on the
$i$-th and $j$-th tensor factors of
$(\cA^!)^{\otimes 3}$.
\end{theorem}

\begin{proof}
The proof proceeds through three steps.

\emph{Step~1: Setup.}
Consider the MC equation on the genus-zero convolution algebra at
degree~$3$ (four marked points: three input + one output at infinity).
The boundary of $\overline{\mathcal{M}}_{0,4}$ has three divisors
$D_s, D_t, D_u$ (Lemma~\ref{V1-lem:thqg-V-m04-structure}).
The MC equation
\begin{equation}% label removed: eq:thqg-V-mc-degree3
D_\cA\Theta_{0,3}
\;+\;
\tfrac{1}{2}\bigl[\Theta_{0,2},\;\Theta_{0,2}\bigr]_{0,3}
\;=\; 0,
\end{equation}
where $[\cdot,\cdot]_{0,3}$ denotes the bracket projected to the
$3$-input component via operadic composition, splits by the collision
filtration into contributions at each boundary stratum.

\emph{Step~2: Arnold computation.}
On $\Conf_3(\C)$, the Arnold relation
(Proposition~\ref{V1-prop:thqg-V-arnold-cohomology}) gives
\begin{equation}% label removed: eq:thqg-V-arnold-3pt
\eta_{12} \wedge \eta_{23}
\;+\;
\eta_{23} \wedge \eta_{31}
\;+\;
\eta_{31} \wedge \eta_{12}
\;=\; 0,
\end{equation}
where $\eta_{ij} = d\log(z_i - z_j)$.
Tensoring with the Koszul-dual Casimir
$\Omega \in (\cA^!)^{\otimes 2}$ on each factor and evaluating
the bracket on $(\cA^!)^{\otimes 3}$, we obtain:
\begin{align}
&\sum_{a,b} \bigl(\Omega_{12}^{ab}\,\Omega_{23}^{cd}
 - \Omega_{23}^{ab}\,\Omega_{12}^{cd}\bigr)\,
 \frac{1}{z_{12}\,z_{23}}
\notag \\
&\qquad +\;
\sum_{a,b} \bigl(\Omega_{23}^{ab}\,\Omega_{31}^{cd}
 - \Omega_{31}^{ab}\,\Omega_{23}^{cd}\bigr)\,
 \frac{1}{z_{23}\,z_{31}}
\notag \\
&\qquad +\;
\sum_{a,b} \bigl(\Omega_{31}^{ab}\,\Omega_{12}^{cd}
 - \Omega_{12}^{ab}\,\Omega_{31}^{cd}\bigr)\,
 \frac{1}{z_{31}\,z_{12}}
\;=\; 0.
% label removed: eq:thqg-V-arnold-tensored
\end{align}

\emph{Step~3: Identification with CYBE.}
Each term in~\eqref{V1-eq:thqg-V-arnold-tensored} is a Lie bracket of
$r$-matrix components:
\begin{align}
\frac{\Omega_{12}\,\Omega_{23}
 - \Omega_{23}\,\Omega_{12}}{z_{12}\,z_{23}}
&\;=\;
\bigl[r_{12}(z_{12}),\;r_{23}(z_{23})\bigr],
% label removed: eq:thqg-V-cybe-term1 \\
\frac{\Omega_{23}\,\Omega_{31}
 - \Omega_{31}\,\Omega_{23}}{z_{23}\,z_{31}}
&\;=\;
\bigl[r_{23}(z_{23}),\;r_{31}(z_{31})\bigr],
% label removed: eq:thqg-V-cybe-term2 \\
\frac{\Omega_{31}\,\Omega_{12}
 - \Omega_{12}\,\Omega_{31}}{z_{31}\,z_{12}}
&\;=\;
\bigl[r_{31}(z_{31}),\;r_{12}(z_{12})\bigr],
% label removed: eq:thqg-V-cybe-term3
\end{align}
where $r_{ij}(z) = \Omega_{ij}/z$ is the leading-order $r$-matrix.
(For algebras with higher-order poles in the $r$-matrix, such as
$\mathrm{Vir}_c$, additional terms appear at higher collision depth;
the leading-order CYBE is the universal statement valid for all
modular Koszul algebras.)

The spectral parameter $z_{ij} = z_i - z_j$ satisfies
$z_{ji} = -z_{ij}$. Skew-symmetry of~$r$ gives
$r(z_{ji}) = r(-z_{ij}) = -P \cdot r(z_{ij}) \cdot P$, where
$P$ is the permutation operator exchanging the two tensor factors.
In particular $z_{31} = -z_{13}$, and
$r_{31}(z_{31}) = r_{31}(-z_{13}) = -r_{13}(z_{13})$
by the skew-symmetry property
(Theorem~\ref{V1-thm:thqg-V-collision-twisting}(ii)).
With this identification, the sum
\eqref{V1-eq:thqg-V-cybe-term1}--\eqref{V1-eq:thqg-V-cybe-term3}
gives exactly the CYBE~\eqref{V1-eq:thqg-V-cybe}.
\end{proof}

\begin{remark}[Gravitational interpretation of the CYBE]
% label removed: rem:thqg-V-cybe-gravitational
\index{classical Yang--Baxter equation!gravitational interpretation}
The CYBE has a gravitational reading.
Consider three line operators at positions $z_1, z_2, z_3$ on the
boundary of the HT theory. In that reading, the binary kernel
$r_{ij}(z_{ij})$ is read as the pairwise graviton-exchange
contribution between the $i$-th and $j$-th insertions. The
CYBE~\eqref{V1-eq:thqg-V-cybe} then states that the three possible
pairwise binary contributions are \emph{consistent}: there is no
tree-level inconsistency in assembling the ternary scattering data
from this binary kernel.

In the celestial/gravitational reading, this same compatibility is
read as a Ward-identity manifestation of the leading celestial soft
factor (Theorem~\ref{thm:thqg-VI-leading-soft}), rather than as an
independent theorem-level statement about tree-level graviton
scattering.
\end{remark}


\subsubsection{The infinitesimal braid relation}

\begin{corollary}[Infinitesimal braid relation;
\ClaimStatusProvedHere]
% label removed: cor:thqg-V-ibr
\index{infinitesimal braid relation}
For any modular Koszul chiral algebra~$\cA$ with Casimir
$\Omega_{ij}$ acting on the $i$-th and $j$-th tensor factors of
$(\cA^!)^{\otimes n}$, the infinitesimal braid relation
\textup{(}IBR\textup{)} holds:
\begin{equation}% label removed: eq:thqg-V-ibr
[\Omega_{ij},\;\Omega_{ik} + \Omega_{jk}] \;=\; 0
\qquad
\text{for all distinct } i,j,k \in \{1,\dotsc,n\}.
\end{equation}
This is a formal consequence of the CYBE at the Casimir level:
setting $z_{12} = z_{13} = z$ (equal spectral parameters) in
\eqref{V1-eq:thqg-V-cybe} and extracting the $z^{-2}$ coefficient
gives the IBR\@.
\end{corollary}

\begin{proof}
The CYBE~\eqref{V1-eq:thqg-V-cybe} with
$r_{ij}(z) = \Omega_{ij}/z$ becomes
\[
\frac{[\Omega_{12},\Omega_{13}]}{z_{12}\,z_{13}}
+
\frac{[\Omega_{12},\Omega_{23}]}{z_{12}\,z_{23}}
+
\frac{[\Omega_{13},\Omega_{23}]}{z_{13}\,z_{23}}
= 0.
\]
Multiply through by $z_{12}\,z_{13}\,z_{23}$ and set
$z_{23} = z_{13} - z_{12}$:
\[
[\Omega_{12},\Omega_{13}]\,z_{23}
+ [\Omega_{12},\Omega_{23}]\,z_{13}
+ [\Omega_{13},\Omega_{23}]\,z_{12}
= 0.
\]
The coefficient of $z_{12}^0\,z_{13}^1$ (equivalently,
taking $z_{12} \to 0$) gives
$[\Omega_{12},\,\Omega_{13} + \Omega_{23}] = 0$,
which is the IBR\@.
\end{proof}

\begin{proposition}[$A_\infty$-enhancement of the CYBE;
\ClaimStatusProvedHere]
% label removed: prop:thqg-V-ainfty-ybe
\index{Yang--Baxter equation!A-infinity enhancement@$A_\infty$ enhancement}
For chiral algebras of shadow class~$C$ or~$M$, the classical
CYBE receives $A_\infty$ corrections from higher collision depths.
The corrected equation at depth~$d$ is
\begin{equation}% label removed: eq:thqg-V-ainfty-ybe
\sum_{k=0}^{d}
\sum_{\substack{i<j<\ell \\
 d_1 + d_2 + d_3 = d \\
 d_1, d_2, d_3 \geq 0}}
\bigl[
 r^{(d_1)}_{ij}(z_{ij}),\;
 r^{(d_2)}_{j\ell}(z_{j\ell})
\bigr]^{(d_3)}
\;=\; 0,
\end{equation}
where $r^{(d)}$ denotes the $r$-matrix at collision depth~$d$ and
$[\cdot,\cdot]^{(d)}$ is the $A_\infty$ bracket at depth~$d$.

For class~$G$ and~$L$: $r^{(d)} = 0$ for $d \geq 1$, and the
$A_\infty$ CYBE reduces to the classical CYBE\@.
For class~$C$: $r^{(1)} \neq 0$ at the quartic level, and the
first correction is the quartic contact term.
For class~$M$: all $r^{(d)}$ are nonzero, and the $A_\infty$ CYBE
is an infinite tower of higher-depth algebraic identities. In the
gravitational reading, this tower is read as organizing the full
nonlinear scattering hierarchy.
\end{proposition}

\begin{proof}
The $A_\infty$ enhancement arises from the full MC equation at all
collision depths, not just depth~$2$. At depth $d$, the MC equation
restricted to $\overline{\mathcal{M}}_{0,4}$ involves compositions
of operations at various intermediate depths summing to~$d$. The
resulting identity is the $A_\infty$ CYBE at total depth~$d$.

For class~$G$: the shadow obstruction tower terminates at degree~$2$
(Corollary~\ref{V1-cor:heisenberg-postnikov-termination}), so all
higher-depth $r$-matrices vanish.
For class~$L$: the shadow obstruction tower terminates at degree~$3$
(Corollary~\ref{V1-cor:affine-postnikov-termination}), so $r^{(d)} = 0$
for $d \geq 2$ and the first potential correction (at degree~$3$,
depth~$1$) is killed by the Jacobi identity.
For class~$C$: the shadow obstruction tower terminates at degree~$4$, and the
quartic contact invariant $\mathfrak{Q}$ provides the first nonzero
$A_\infty$ correction at total depth~$2$.
For class~$M$: the tower is infinite
(Theorem~\ref{V1-thm:nms-virasoro-quintic-forced}).
\end{proof}


\subsubsection{Higher-degree CYBE}

\begin{definition}[$n$-point Yang--Baxter identity]
% label removed: def:thqg-V-n-point-ybe
\index{Yang--Baxter equation!n-point@$n$-point}
For $n \geq 3$ points on $\overline{\mathcal{M}}_{0,n+1}$, the
\emph{$n$-point Yang--Baxter identity} is the MC equation projected
to collision depth $n - 1$:
\begin{equation}% label removed: eq:thqg-V-n-point-ybe
\sum_{\substack{T \in \mathcal{T}_n \\
 |E_{\mathrm{int}}(T)| = n-2}}
\frac{1}{|\Aut(T)|}
\prod_{e \in E_{\mathrm{int}}(T)}
r_{e_-,e_+}(z_{e_-,e_+})
\;=\; 0,
\end{equation}
where $\mathcal{T}_n$ is the set of binary rooted trees with $n$
leaves, $e_-$ and $e_+$ are the endpoints of the internal edge~$e$,
and the product is over internal edges.
At $n = 3$: this is the CYBE (three binary trees, one internal edge
each).
At $n = 4$: this gives the \emph{quartic YBE}, involving the
five binary trees with two internal edges.
\end{definition}

\begin{proposition}[$n$-point YBE from boundary of
$\overline{\mathcal{M}}_{0,n+1}$; \ClaimStatusProvedHere]
% label removed: prop:thqg-V-n-point-ybe-proof
\index{Yang--Baxter equation!n-point proof@$n$-point proof}
The $n$-point YBE is a consequence of the vanishing
\begin{equation}% label removed: eq:thqg-V-boundary-sum
\sum_{D \in \partial\overline{\mathcal{M}}_{0,n+1}}
\int_D
\prod_{e}
r_{e_-,e_+}(z_{e_-,e_+})
\;=\; 0,
\end{equation}
which follows from $\partial^2 = 0$ on the chain level
\textup{(}the boundary of a boundary is empty\textup{)}, together
with the MC equation at all intermediate degrees.
\end{proposition}

\begin{proof}
The MC equation $D\Theta + \tfrac{1}{2}[\Theta,\Theta] = 0$ at
degree~$n$, projected to the maximal collision depth, gives a sum over
all maximally degenerate boundary strata of
$\overline{\mathcal{M}}_{0,n+1}$. These strata are indexed by binary
trees with $n$ leaves. The identity $\partial^2 = 0$ on
$C_*(\overline{\mathcal{M}}_{0,n+1})$ ensures that the total
boundary contribution vanishes. Each boundary stratum contributes a
product of $r$-matrices along internal edges, giving
\eqref{V1-eq:thqg-V-boundary-sum}.

The identification of the boundary sum with the combinatorial formula
\eqref{V1-eq:thqg-V-n-point-ybe} follows from the factorization of
boundary strata: each boundary divisor of $\overline{\mathcal{M}}_{0,n+1}$
is a product of lower-dimensional moduli spaces, and the $r$-matrix
contributions factor accordingly.
\end{proof}


% ======================================================================
%
% SUBSECTION 4: THE MODULAR DG-SHIFTED YANGIAN
%
% ======================================================================

\subsection{The modular dg-shifted Yangian}
% label removed: subsec:thqg-V-modular-yangian
\index{Yangian!modular dg-shifted|textbf}
\index{Yangian!modular higher-genus completion}

The genus-zero $r$-matrix is the leading term of a richer object:
the \emph{modular dg-shifted Yangian} $\Ymod_\cA$, an ambient
pronilpotent dg Lie algebra that organizes the higher-genus binary
collision data. The formal series of genus-refined binary
projections is defined on this surface; identifying it with a genuine
modular kernel and proving the stable-graph Yang--Baxter equation are
additional frontier inputs. The construction assembles the
collision-filtration data of
\S\ref{V1-subsec:thqg-V-collision-filtration} into a single algebraic
object.

\subsubsection{Pronilpotent completion}

\begin{definition}[Modular dg-shifted Yangian for $\cA$;
\ClaimStatusProvedHere]
% label removed: def:thqg-V-modular-yangian
\index{Yangian!modular dg-shifted!definition}
Let $\cA$ be a modular Koszul chiral algebra.
The \emph{modular dg-shifted Yangian} of~$\cA$ is the pronilpotent
dg Lie algebra
\begin{equation}% label removed: eq:thqg-V-modular-yangian
\Ymod_\cA
\;:=\;
\varprojlim_N\;
\mathcal{L}_\cA^{\mathrm{mod}} \big/ F^{N+1}\!
 \mathcal{L}_\cA^{\mathrm{mod}},
\end{equation}
where
\begin{equation}% label removed: eq:thqg-V-convolution
\mathcal{L}_\cA^{\mathrm{mod}}
\;:=\;
\prod_{2g-2+n > 0}
\Hom\!\bigl(
 C_\bullet(\overline{\mathcal{M}}_{g,n+1})
 \otimes (\cA^!)^{\widehat{\otimes}\,n},\;
 \cA^!
\bigr)[1]
\end{equation}
is the modular convolution dg Lie algebra
\textup{(}Definition~\textup{\ref*{V1-def:modular-convolution-dg-lie})}
with the filtration
$F^N \mathcal{L}_\cA^{\mathrm{mod}}
:= \prod_{2g-2+n \geq N}(\cdots)$.

The genus-zero truncation is the
(non-modular) dg-shifted Yangian:
\begin{equation}% label removed: eq:thqg-V-genus-zero-truncation
\Ydg_\cA
\;=\;
\Ymod_\cA\big|_{g=0}
\;=\;
\prod_{n \geq 2}
\Hom_{\Sigma_n}\!\bigl(
 C_\bullet(\overline{\mathcal{M}}_{0,n+1})
 \otimes (\cA^!)^{\widehat{\otimes}\,n},\;
 \cA^!
\bigr)[1].
\end{equation}
\end{definition}

\begin{remark}[Comparison with
Definition~\textup{\ref*{V1-def:modular-yangian-pro}}]
% label removed: rem:thqg-V-yangian-comparison
Definition~\ref{V1-def:thqg-V-modular-yangian} is a restatement
of the ambient pronilpotent dg Lie package in
Definition~\ref*{V1-def:modular-yangian-pro} in the gravitational
context. The key structural point is that the Lie filtration
$F^{N_1} \times F^{N_2} \to F^{N_1 + N_2}$
(Lemma~\ref{V1-lem:thqg-V-collision-lie-compatibility}) ensures the
pronilpotent completion is well-defined. What is \emph{not} proved
here is that the completion already comes equipped with a modular
kernel satisfying an MC equation; that stronger package remains on the
conjectural modular-Yangian surface.
\end{remark}

\begin{remark}[Bialgebra structure: no antipode, and the name ``Yangian'']
% label removed: rem:thqg-V-yangian-no-antipode
\index{Yangian!no antipode}
\index{antipode!absence in dg-shifted Yangian}
Two terminological points require clarification.

\smallskip\noindent
\emph{(i) Bialgebra, not Hopf algebra.}
The dg-shifted Yangian $\Ydg_\cA$ carries a coproduct
$\Delta_z$ (from the operadic composition on the
convolution algebra) and a counit $\varepsilon$
(from the augmentation), making it a dg bialgebra.
It does \emph{not} carry an antipode in general: the
Drinfeld Yangian $Y(\fg)$ has an antipode $S$ by the
finiteness of its Drinfeld generators, but the $A_\infty$
analogue for arbitrary modular Koszul algebras is not
known to admit one (cf.\
Remark~\ref{rem:dg-yangian-antipode}). For the
representation-theoretic applications in this chapter
(line-operator modules, spectral braiding, the RTT
relation), the bialgebra structure is sufficient: the
antipode would be needed for dualization of modules
and for constructing a ribbon structure, neither of
which is required here.

\smallskip\noindent
\emph{(ii) The name ``Yangian.''}
We use ``dg-shifted Yangian'' to emphasize the structural
parallel with the Drinfeld Yangian $Y(\fg)$: both are
pronilpotent completions of convolution algebras. For
$\cA = \widehat{\fg}_k$ at generic level,
$\Ydg_\cA$ recovers the standard Yangian $Y(\fg)$
(Computation~\ref{V1-comp:thqg-V-affine-yangian}). For
general~$\cA$, the object is not an ``affine Yangian''
in the sense of Guay--Nakajima--Wendlandt (which is a
specific algebra associated to $\widehat{\fg}$ with its
loop parameter); rather, it is the genus-zero binary
projection of the modular convolution algebra. The
modular extension $\Ymod_\cA$, which includes all genera,
has no classical precedent. The name ``gravitational
Yangian'' for the physical incarnation is adopted to
distinguish it from both the classical Yangian and
the affine Yangian.

\smallskip\noindent
\emph{(iii) Consistency with
Definition~\textup{\ref{def:dg_Yangian}}.}
The genus-zero truncation $\Ydg_\cA$ satisfies the axioms of
Definition~\ref{def:dg_Yangian}: the translation automorphism
$\tau_z$ is the spectral-parameter shift on the $r$-matrix,
the coproduct arises from operadic composition, and the YBE is
Theorem~\ref{V1-thm:thqg-V-cybe-from-arnold}. The modular
extension $\Ymod_\cA$ is a genuine enlargement: the higher-genus
components are not part of the abstract dg-shifted Yangian
structure of Definition~\ref{def:dg_Yangian} but are dictated
by the MC equation in the full modular convolution algebra.
\end{remark}

\begin{proposition}[Differential structure;
\ClaimStatusProvedHere]
% label removed: prop:thqg-V-yangian-differential
\index{Yangian!modular dg-shifted!differential}
The differential on $\Ymod_\cA$ has two summands:
\begin{equation}% label removed: eq:thqg-V-yangian-diff
d_{\Ymod}
\;=\;
d_{\mathrm{int}} + d_{\mathrm{mod}},
\end{equation}
where:
\begin{enumerate}[label=\textup{(\roman*)}]
\item $d_{\mathrm{int}}(\phi)$ acts by the $A_\infty$ differential
 on $\cA^!$ applied to the inputs and output of~$\phi$;
\item $d_{\mathrm{mod}}(\phi)$ acts by precomposition with the
 boundary operator
 $\partial\colon C_*(\overline{\mathcal{M}}_{g,n+1})
 \to \coprod C_*(\overline{\mathcal{M}}_{g_1,n_1+1}
 \times \overline{\mathcal{M}}_{g_2,n_2+1})$.
\end{enumerate}
The identity $d_{\Ymod}^2 = 0$ is equivalent to the $A_\infty$
structure equations on $\cA^!$ together with $\partial^2 = 0$ on the
chain complexes of moduli spaces.
\end{proposition}

\begin{proof}
$d_{\mathrm{int}}^2 = 0$ because $\cA^!$ is an $A_\infty$-algebra.
$d_{\mathrm{mod}}^2 = 0$ because $\partial^2 = 0$ on chains.
The cross-term
$d_{\mathrm{int}} \circ d_{\mathrm{mod}} +
d_{\mathrm{mod}} \circ d_{\mathrm{int}} = 0$
holds because the $A_\infty$ structure on $\cA^!$ is compatible
with the modular operad structure: the boundary operator on moduli
spaces composes with the $A_\infty$ multiplication maps. This
compatibility is the content of the bar complex being a
module over the Feynman transform of the modular operad
(Theorem~\ref{V1-thm:bar-modular-operad}).
\end{proof}


\subsubsection{The formal genus-refined binary shadow series}

\begin{theorem}[Formal genus-refined binary shadow series
 $R^{\mathrm{bin}}_\cA$;
\ClaimStatusProvedHere]
% label removed: thm:thqg-V-pro-mc-element
\index{Yangian!binary shadow series|textbf}
The universal MC element $\Theta_\cA \in \MC(\gAmod)$
\textup{(}Theorem~\textup{\ref*{V1-thm:mc2-bar-intrinsic})}
determines formal genus-refined binary coefficients
$r_{\cA,g}(z)$ and hence a formal series
\begin{equation}% label removed: eq:thqg-V-pro-mc
R^{\mathrm{bin}}_\cA(z;\hbar)
\;:=\;
\sum_{g \geq 0}
\hbar^{2g}\,r_{\cA,g}(z)
\;\in\;
\Ymod_\cA[\![\hbar^2]\!],
\end{equation}
where $r_{\cA,g}(z)$ is the genus-$g$ component of
$\Theta_\cA$ projected to the binary collision stratum.
\end{theorem}

\begin{proof}
The bar-intrinsic MC element $\Theta_\cA = D_\cA - d_0$
satisfies $D_\cA\Theta_\cA + \tfrac{1}{2}[\Theta_\cA,\Theta_\cA] = 0$
in $\gAmod$. Projecting $\Theta_\cA$ to the genus-$g$ binary
collision stratum produces the coefficient $r_{\cA,g}(z)$.
Because $\Ymod_\cA$ is the pronilpotent completion of the modular
convolution dg Lie algebra, these coefficients assemble into the
formal $\hbar$-adic series~\eqref{V1-eq:thqg-V-pro-mc}.
\end{proof}

\begin{corollary}[Genus-$0$ component = classical $r$-matrix;
\ClaimStatusProvedHere]
% label removed: cor:thqg-V-genus-zero-r
\index{r-matrix@$r$-matrix!genus-zero component}
The genus-$0$ component $r_{\cA,0}(z)$ is the classical $r$-matrix
of Theorem~\ref{V1-thm:thqg-V-collision-twisting}:
$r_{\cA,0}(z) = r_\cA(z)$.
\end{corollary}

\begin{proof}
Restrict $R^{\mathrm{bin}}_\cA$ to $\hbar = 0$.
\end{proof}

\begin{remark}[Heuristic genus-$1$ correction;
\ClaimStatusHeuristic]
% label removed: cor:thqg-V-genus-one-correction
\index{r-matrix@$r$-matrix!genus-one correction}
The first higher-genus binary coefficient is expected to involve the
curvature $\kappa(\cA)$, the genus-$1$ Hessian
$\delta H^{(1)}_\cA$, and quartic shadow corrections in the form
\begin{equation}% label removed: eq:thqg-V-genus-one-r
r_{\cA,1}(z)
\;=\;
\kappa(\cA) \cdot r_{\cA,0}(z)
\;+\;
\delta H^{(1)}_\cA \cdot \partial_z r_{\cA,0}(z)
\;+\;
(\text{quartic corrections}).
\end{equation}
For class~$G$ and~$L$ algebras, where the quartic corrections vanish
on the relevant binary-shadow surface, this predicts a controlled
deformation of the classical genus-zero $r$-matrix. What is not
proved here is that these coefficients assemble into a genuine
modular Maurer--Cartan kernel.
\end{remark}


\subsubsection{The stable-graph Yang--Baxter equation}

\begin{definition}[Stable-graph YBE]
% label removed: def:thqg-V-stable-graph-ybe
\index{Yang--Baxter equation!stable-graph}
On the conjectural modular-Yangian surface, the
\emph{stable-graph Yang--Baxter equation} (sgYBE) is the MC
equation for a modular kernel
$R^{\mathrm{mod}}_\cA(z;\hbar)\in\MC(\Ymod_\cA)$ expanded by genus:
at genus~$g$ and degree~$n$, it reads
\begin{equation}% label removed: eq:thqg-V-sgybe-expanded
d_{\mathrm{int}}\,r_{\cA,g}^{(n)}
\;+\;
\sum_{\Gamma \in \mathcal{G}^{\mathrm{st}}_{g,n}}
\frac{1}{|\Aut(\Gamma)|}
\prod_{v \in V(\Gamma)} r_{\cA,g_v}^{(n_v)}
\prod_{e \in E(\Gamma)} P_e
\;=\; 0,
\end{equation}
where $\mathcal{G}^{\mathrm{st}}_{g,n}$ is the set of stable
graphs of genus~$g$ with $n$ legs, the vertex labels $r_{\cA,g_v}^{(n_v)}$
are the pro-MC components, and $P_e$ is the edge propagator.

At genus~$0$, degree~$3$: the sgYBE would reduce to the classical
CYBE (three stable graphs, each a single edge connecting two
trivalent vertices).

At genus~$1$, degree~$2$: the sgYBE would give the
\emph{quantum Yang--Baxter equation} (qYBE), encoding the
one-loop correction to the $r$-matrix. The single stable graph
is a loop (one vertex, one self-edge), and the equation reads
\begin{equation}% label removed: eq:thqg-V-qybe
d_{\mathrm{int}}\,r_{\cA,1}^{(2)}
\;+\;
\tfrac{1}{2}\,\Tr_{12}\,
\bigl[r_{\cA,0}^{(2)},\,r_{\cA,0}^{(2)}\bigr]
\;+\;
\kappa(\cA) \cdot r_{\cA,0}^{(2)}
\;=\; 0,
\end{equation}
where $\Tr_{12}$ traces over the loop-contracted tensor factors and
the $\kappa$-term is the curvature contribution from the degenerating
cycle.
\end{definition}

\begin{remark}[The sgYBE and its gravitational-reading organization]
% label removed: rem:thqg-V-sgybe-scattering
\index{gravitational scattering!stable-graph YBE}
If a modular kernel satisfying the sgYBE is supplied, then it
organizes the corresponding candidate gravitational amplitudes by
topology. At genus~$0$: tree-level graviton exchange, controlled by
the classical $r$-matrix. At genus~$1$: one-loop graviton
self-energy and vertex corrections, controlled by $r_{\cA,1}$.
At genus~$g$: the $g$-loop gravitational amplitude, with the
combinatorics of stable graphs providing the Feynman-diagrammatic
structure. The present chapter does not prove the existence of such a
kernel; it only isolates the ambient modular dg Lie package in which
that higher-genus completion would live.
\end{remark}

\begin{proposition}[Formal \texorpdfstring{$\hbar$}{hbar}-adic
 convergence of the candidate binary shadow series;
\ClaimStatusProvedHere]
% label removed: prop:thqg-V-pro-mc-convergence
\index{Yangian!binary shadow series!convergence}
The formal series $R^{\mathrm{bin}}_\cA(z;\hbar)$ converges in
the $\hbar$-adic topology of $\Ymod_\cA$: at each genus~$g$,
the component $r_{\cA,g}(z)$ is a well-defined element of the
finite-dimensional $\Hom$ space, and the genus series
$\sum_g \hbar^{2g}\,r_{\cA,g}(z)$ is a formal power series in
$\C[\![\hbar^2]\!]$ with coefficients in the genus-zero Yangian.

For the Heisenberg algebra: $r_{\mathcal{H},g}(z) = 0$ for
$g \geq 1$ \textup{(}the genus tower terminates at genus~$0$
for the $r$-matrix\textup{)}.

For affine Kac--Moody at generic level: $r_{\widehat{\fg}_k,g}(z)$
is polynomial in $1/(k+h^\vee)$ at each genus, with degree
controlled by the graph-sum formula
\textup{(}Theorem~\textup{\ref{V1-thm:recursive-existence})}.
\end{proposition}

\begin{proof}
The convergence in the $\hbar$-adic topology is automatic from the
filtration: $r_{\cA,g}$ lies in $F^{2g-2+n}\,\Ymod_\cA$, and the
filtration is exhaustive and separated. The Heisenberg vanishing
for $g \geq 1$ follows from the Gaussian nature of the Heisenberg:
the genus-$g$ sewing amplitude is a determinant
(Theorem~\ref{thm:heisenberg-sewing}), which contributes to the
scalar shadow $\kappa$ but not to the $r$-matrix component.

For affine KM, the genus-$g$ component is computed by the graph-sum
formula with $|V_{\mathrm{int}}| = 2g - 2 + n$ vertices; each
vertex carries a factor of the Casimir $\Omega/(k+h^\vee)$, giving
the polynomial dependence on $1/(k+h^\vee)$.
\end{proof}


\subsubsection{Weight filtration and Yangian generators}

\begin{construction}[Yangian generators from the weight filtration;
\ClaimStatusProvedHere]
% label removed: constr:thqg-V-yangian-generators
\index{Yangian!generators from weight filtration}
The weight filtration on $\Ydg_\cA$ at genus zero
(the degree filtration) produces Yangian generators.
Let $\{e_\alpha\}$ be a basis of $\cA^!$ homogeneous with respect
to the conformal weight grading.

\emph{Level~$0$ generators.}
$T_{\alpha\beta}^{(0)} := \langle e_\alpha,\,r_\cA(z)\,e_\beta\rangle
= \Omega_{\alpha\beta}/z$.
These are the $\fg$-currents at degree~$2$.

\emph{Level~$r$ generators.}
$T_{\alpha\beta}^{(r)}$ is the coefficient of $z^{-r-1}$ in the
Laurent expansion of $r_\cA(z)$ at $z = 0$:
\begin{equation}% label removed: eq:thqg-V-yangian-generators
r_\cA(z)
\;=\;
\sum_{r \geq 0}
T^{(r)}\,z^{-r-1},
\qquad
T^{(r)} \in (\cA^!)^{\otimes 2}.
\end{equation}
The commutation relations of the $T^{(r)}$ are determined by the
CYBE and its higher-degree analogues. For $\cA = \widehat{\fg}_k$:
$T^{(0)} = \Omega/(k+h^\vee)$ and $T^{(r)} = 0$ for $r \geq 1$,
so the Yangian is generated by the level-$0$ Casimir alone.
For $\cA = Y(\fg)$: all $T^{(r)}$ are nonzero, recovering the
standard Yangian generators
(Definition~\ref{V1-def:yangian}).
\end{construction}

\begin{proposition}[RTT relation from the genus-zero collision YBE;
\ClaimStatusProvedHere]
% label removed: prop:thqg-V-rtt-from-sgybe
\index{RTT relation!from genus-zero YBE}
The genus-zero collision/Yang--Baxter package, expressed in terms
of the transfer matrix
$T_\cA(z) = \sum_r T^{(r)} z^{-r}$, gives the RTT relation:
\begin{equation}% label removed: eq:thqg-V-rtt
r_{12}(z_{12})\,T_1(z_1)\,T_2(z_2)
\;-\;
T_2(z_2)\,T_1(z_1)\,r_{12}(z_{12})
\;=\;
T_1(z_1)\,r_{12}(z_{12})\,T_2(z_2)
\;-\;
T_2(z_2)\,r_{12}(z_{12})\,T_1(z_1),
\end{equation}
which rearranges to the standard RTT form
$r_{12}(z_{12})\,T_1(z_1)\,T_2(z_2)
= T_2(z_2)\,T_1(z_1)\,r_{12}(z_{12})$
when $r$ is the fundamental $R$-matrix.
\end{proposition}

\begin{proof}
The CYBE at genus~$0$ involves $r_{12}$, $r_{13}$, and $r_{23}$ with
$r_{ij}(z) = \Rescoll_{0,2}(\Theta_\cA)|_{ij}$. When acting on a
module via the evaluation representation, $r_{i3}(z)$ becomes the
action of $T(z)$ on the $i$-th factor. The three-term CYBE then
becomes the RTT relation after distributing the Lie bracket over the
tensor product. The expansion in modes $z^{-r-1}$ gives the
quadratic RTT relations among the generators $T^{(r)}$
(cf.~Definition~\ref{V1-def:yangian-rtt}).
\end{proof}


% ======================================================================
%
% SUBSECTION 5: GRAVITATIONAL YANGIAN FOR VIRASORO
%
% ======================================================================

\subsection{Gravitational Yangian for Virasoro}
% label removed: subsec:thqg-V-virasoro-yangian
\index{Virasoro!gravitational Yangian|textbf}
\index{gravitational Yangian!Virasoro}

The Virasoro algebra is the simplest class-$M$ algebra: its shadow
tower is infinite, the $r$-matrix has higher-order poles, and the
genus-zero gravitational Yangian exhibits the richest explicit binary
collision data among the standard test cases. This subsection
computes the explicit Virasoro binary kernel, the ternary shadow,
and the expected representation-theoretic model.


\subsubsection{The Virasoro binary kernel and its gravitational reading}

\begin{construction}[Explicit Virasoro binary kernel;
\ClaimStatusProvedHere]
\label{constr:thqg-V-vir-r-explicit}
\index{Virasoro!r-matrix@$r$-matrix!explicit}
\index{r-matrix@$r$-matrix!Virasoro!explicit construction}
For $\mathrm{Vir}_c$, the Koszul dual is $\mathrm{Vir}_{26-c}$
with Virasoro generators $\{L_n\}_{n \in \Z}$ satisfying
$[L_m, L_n] = (m-n)\,L_{m+n}
+ \frac{26-c}{12}\,(m^3-m)\,\delta_{m+n,0}$.

The Virasoro binary kernel $\rgrav_{\mathrm{Vir}}(z)$ acts on
$\mathrm{Vir}_{26-c} \otimes \mathrm{Vir}_{26-c}$ and is determined
by the collision residue of $\Theta_{\mathrm{Vir}}$:
\begin{equation}% label removed: eq:thqg-V-vir-r-explicit
\rgrav_{\mathrm{Vir}}(z)
\;=\;
\sum_{m \geq 0}
\frac{1}{z^{m+1}}\,
\Omega_m,
\end{equation}
where the $\Omega_m$ are the Virasoro Casimir components:
\begin{align}
\Omega_0
&\;=\;
\sum_{n \in \Z}
L_{-n} \otimes L_n
\;\in\;
(\mathrm{Vir}_{26-c})^{\widehat{\otimes}\,2},
% label removed: eq:thqg-V-vir-omega0 \\
\Omega_1
&\;=\;
\sum_{n \in \Z}
n\,L_{-n} \otimes L_n,
% label removed: eq:thqg-V-vir-omega1 \\
\Omega_m
&\;=\;
\sum_{n \in \Z}
\binom{n+m-1}{m}\,L_{-n-m} \otimes L_n
\qquad (m \geq 0).
% label removed: eq:thqg-V-vir-omegam
\end{align}
The $m = 0$ term is the leading Casimir pole; the $m \geq 1$ terms
are the higher poles arising from the Virasoro fourth-order OPE
singularity.

The $r$-matrix satisfies:
\begin{enumerate}[label=\textup{(\roman*)}]
\item \emph{Residue.}\;
 $\Res_{z=0}\,\rgrav_{\mathrm{Vir}}(z) = \Omega_0$, the
 Virasoro Casimir.
\item \emph{Skew-symmetry.}\;
 $\rgrav_{\mathrm{Vir}}(z)_{12} = -\rgrav_{\mathrm{Vir}}(-z)_{21}$.
\item \emph{CYBE.}\;
 The classical Yang--Baxter equation holds:
 $[\rgrav_{12}(z_{12}),\,\rgrav_{13}(z_{13})]
 + [\rgrav_{12}(z_{12}),\,\rgrav_{23}(z_{23})]
 + [\rgrav_{13}(z_{13}),\,\rgrav_{23}(z_{23})]
 = 0.$
\end{enumerate}
\end{construction}

In the gravitational reading, the kernel
$\rgrav_{\mathrm{Vir}}(z)$ is read as the binary graviton-exchange
kernel for the Virasoro HT package.

\begin{proof}[Verification of the CYBE]
The CYBE for $\rgrav_{\mathrm{Vir}}$ follows from
Theorem~\ref{V1-thm:thqg-V-cybe-from-arnold} applied to the Virasoro
OPE. Explicitly, the three terms of the CYBE evaluate to:
\begin{align}
&[\rgrav_{12}(z_{12}),\;\rgrav_{13}(z_{13})]
\;=\;
\sum_{m_1,m_2}
\frac{[\Omega_{m_1,12},\,\Omega_{m_2,13}]}
 {z_{12}^{m_1+1}\,z_{13}^{m_2+1}},
% label removed: eq:thqg-V-vir-cybe-1
\end{align}
and similarly for the other two terms. Each bracket
$[\Omega_{m_1},\,\Omega_{m_2}]$ is computed using the Virasoro
commutation relations; the resulting Laurent series in $z_{12}, z_{13}$
satisfies the partial-fraction identity
\[
\frac{1}{z_{12}^{a+1}\,z_{13}^{b+1}}
+
\frac{1}{z_{12}^{a+1}\,z_{23}^{b+1}}
+
\frac{1}{z_{13}^{a+1}\,z_{23}^{b+1}}
=
\text{(lower-order poles)},
\]
and the cancellation at each order follows from the Arnold relation
on $\overline{\mathcal{M}}_{0,4}$.

The key identity is the \emph{Virasoro partial-fraction relation}:
for the Casimir components,
\begin{equation}% label removed: eq:thqg-V-vir-partial-fraction
\sum_{\text{cyclic}(1,2,3)}
[\Omega_{m,12},\,\Omega_{m',13}]
\;=\;
0
\end{equation}
when $m + m' = $ const, which is the Virasoro
specialization of the IBR
(Corollary~\ref{V1-cor:thqg-V-ibr}).
\end{proof}


\subsubsection{The ternary shadow and its gravitational reading}

\begin{computation}[Ternary shadow and its gravitational reading;
\ClaimStatusProvedHere]
% label removed: comp:thqg-V-three-graviton
\index{graviton!three-graviton vertex}
\index{Virasoro!three-graviton vertex}
The ternary component of
$\Theta_{\mathrm{Vir}}$ at genus~$0$, extracted via the collision
filtration at depth~$2$. On the space of three Virasoro
primaries of conformal weights $h_1, h_2, h_3$, this operator is:
\begin{equation}% label removed: eq:thqg-V-three-graviton
V_3(h_1, h_2, h_3;\, z_1, z_2, z_3)
\;=\;
\sum_{i<j}
\frac{h_k}{(z_i - z_j)^2}
\;+\;
\frac{1}{z_i - z_j}\,\partial_{z_j},
\end{equation}
where $\{i,j,k\} = \{1,2,3\}$. This is the standard Virasoro
three-point function operator, derived from the shadow connection
$\nabla^{\mathrm{hol}}_{0,3}$
(by the Volume~I Virasoro/BPZ comparison proposition).

The corresponding flat section for primaries $\phi_{h_i}$ is:
\begin{equation}% label removed: eq:thqg-V-three-graviton-amplitude
\langle \phi_{h_1}(z_1)\,\phi_{h_2}(z_2)\,\phi_{h_3}(z_3) \rangle
\;=\;
\frac{C_{h_1 h_2 h_3}}
{z_{12}^{h_1+h_2-h_3}\,z_{13}^{h_1+h_3-h_2}\,z_{23}^{h_2+h_3-h_1}},
\end{equation}
where $C_{h_1 h_2 h_3}$ is the OPE coefficient. This is the flat
section of $\nabla^{\mathrm{hol}}_{0,3}$.

\emph{Gravitational interpretation.}
In the gravitational reading, this ternary shadow is read as the
cubic coupling of three gravitons (stress-tensor insertions) on the
boundary of the HT theory. The denominator structure
$z_{ij}^{-h_i - h_j + h_k}$ is then read as the corresponding
holomorphic momentum-space graviton vertex: the conformal weights
play the role of graviton helicities, and the $z_{ij}$ are the
holomorphic momenta.
\end{computation}

\begin{computation}[Quartic contact term and its gravitational reading;
\ClaimStatusProvedHere]
% label removed: comp:thqg-V-quartic-graviton
\index{graviton!quartic vertex}
\index{quartic contact invariant!graviton}
\index{Virasoro!quartic contact invariant!graviton interpretation}
The degree-$4$ component of
$\Theta_{\mathrm{Vir}}$ at genus~$0$, collision depth~$2$.
Beyond the factorized contribution
(product of two ternary shadows), there is a
\emph{contact} contribution that does not factorize through
any channel:
\begin{equation}% label removed: eq:thqg-V-quartic-contact
V_4^{\mathrm{contact}}(z_1, z_2, z_3, z_4)
\;=\;
\frac{10}{c(5c+22)}\,
\Lambda(z_1, z_2, z_3, z_4),
\end{equation}
where $\Lambda$ is the quartic contact form on
$\overline{\mathcal{M}}_{0,5}$ and the coefficient
$Q^{\mathrm{contact}}_{\mathrm{Vir}} = 10/[c(5c+22)]$ is the
quartic contact invariant
(Computation~\ref{V1-comp:holographic-ss-vir}).

In the gravitational reading, this contact term is read as the first
obstruction to the binary gravitational $r$-matrix being a
\emph{complete} description of the scattering: for class-$M$
algebras, the binary $r$-matrix is then supplemented by ternary,
quartic, and higher shadow data. The quartic contact invariant
measures the strength of this obstruction in that reading.

\emph{Divergence at special central charges.}
The denominator $c(5c+22)$ vanishes at $c = 0$ and
$c = -22/5$. At $c = 0$: the Virasoro algebra degenerates and the
quartic contact is ill-defined. At $c = -22/5$:
the $(2,5)$ minimal model, where the null vector at level~$5$ forces
special behavior. Away from these values, the quartic contact is a
smooth function of~$c$.
\end{computation}


\subsubsection{Highest-weight models for the gravitational Yangian}

\begin{definition}[Expected Virasoro highest-weight model]
% label removed: def:thqg-V-gravitational-hw
\index{gravitational Yangian!highest-weight modules}
In the expected Virasoro realization of the line-side package, one is
led to highest-weight modules generated by a vector
$|h\rangle$ satisfying:
\begin{enumerate}[label=\textup{(\roman*)}]
\item $L_n\,|h\rangle = 0$ for $n > 0$;
\item $L_0\,|h\rangle = h\,|h\rangle$;
\item the $r$-matrix acts as
 $\rgrav_{\mathrm{Vir}}(z)\,|h\rangle \otimes |h'\rangle
 = \frac{h + h'}{z}\,|h\rangle \otimes |h'\rangle
 + (\text{lower poles in }z)$.
\end{enumerate}
In the gravitational reading, these highest-weight modules would model
line objects of conformal weight~$h$.
\end{definition}

\begin{remark}[Expected Verma-module model]
% label removed: prop:thqg-V-verma-gravitational
\index{Verma module!gravitational representation}
In the same expected Virasoro realization, the Verma module $M(h,c)$
is the natural model highest-weight object. The binary shadow action
is expected to take an evaluation-type form
\begin{equation}% label removed: eq:thqg-V-verma-evaluation
\rgrav_{\mathrm{Vir}}(z)\big|_{M(h,c) \otimes M(h',c)}
\;=\;
\frac{1}{z}\,\Omega_0\big|_{M(h,c) \otimes M(h',c)}
\;+\;
\text{(higher-order corrections)}.
\end{equation}
For generic~$h$, one expects the resulting model to be faithful on
$M(h,c)$. For degenerate representations ($h = h_{r,s}$, the Kac
table values), the null vectors generate the ideal, and the quotient
module $L(h_{r,s},c)$ is expected to carry the corresponding truncated
binary shadow action. The PBW and null-vector discussion explains why
this is the natural Virasoro model, but the live surface does not
isolate that realization as a separate proved theorem.
\end{remark}

\begin{remark}[Expected Virasoro highest-weight category]
% label removed: rem:thqg-V-gravitational-rep-category
\index{gravitational Yangian!representation category}
In that expected Virasoro realization, the highest-weight category is
controlled by two parameters: the conformal weight $h$ and the central
charge~$c$. At fixed~$c$, the expected structure is:
\begin{enumerate}[label=\textup{(\roman*)}]
\item \emph{Generic $h$.}
 The Verma module $M(h,c)$ would be irreducible and carry the full
 expected highest-weight model.
\item \emph{Degenerate $h = h_{r,s}$.}
 The irreducible quotient $L(h_{r,s},c)$ would carry a truncated
 expected model. The BPZ null-vector differential equations would
 arise from the truncation of the shadow connection
 $\nabla^{\mathrm{hol}}_{0,n}$ to that quotient.
\item \emph{Rational $c$ (minimal models).}
 The category would have finitely many irreducibles, and the fusion
 rules $N_{ij}^k$ would govern the tensor product of these expected
 model objects. In the same expected realization, the gravitational
 $R$-matrix would reduce to a finite-dimensional matrix acting on the
 finite representation ring.
\end{enumerate}
This structure is the Virasoro specialization of the general
line category, modeled here by modules for the open-colour dual
$\cA^!$ in the sense of
Definition~\ref{V1-def:holographic-modular-koszul-datum}.
\end{remark}


\subsubsection{Self-duality at $c = 13$}

\begin{proposition}[Self-duality of the gravitational Yangian at
$c = 13$; \ClaimStatusProvedHere]
% label removed: prop:thqg-V-c13-self-duality
\index{Virasoro!self-duality at c=13@self-duality at $c = 13$}
\index{gravitational Yangian!self-duality}
At $c = 13$, the Koszul dual $\mathrm{Vir}_{26-c} = \mathrm{Vir}_{13}$
is isomorphic to the original: $\mathrm{Vir}_{13}^! \cong
\mathrm{Vir}_{13}$. The gravitational Yangian becomes self-dual:
\begin{equation}% label removed: eq:thqg-V-c13-self-dual
\Ydg_{\mathrm{Vir}_{13}}
\;\cong\;
\Ydg_{\mathrm{Vir}_{13}^!}
\;\cong\;
\Ydg_{\mathrm{Vir}_{13}},
\end{equation}
and the gravitational $r$-matrix satisfies the enhanced symmetry
$\rgrav(z) = -\rgrav(-z)^{21}$ without change of central charge.

The fixed-point equation $c = 26 - c$ has the unique solution
$c = 13$. This is not $c = 26$ (the critical string): the
self-duality is the Koszul self-duality of the Virasoro algebra,
not the vanishing of the total central charge. At $c = 26$, the
Koszul dual is $\mathrm{Vir}_0$, which is degenerate.
\end{proposition}

\begin{proof}
The Koszul dual of $\mathrm{Vir}_c$ is $\mathrm{Vir}_{26-c}$
(Remark~\ref{V1-rem:virasoro-resonance-model}).
Self-duality requires $c = 26 - c$, i.e., $c = 13$.
The dg-shifted Yangian is a functor of the Koszul dual
$\cA^!$, so self-duality of $\cA^!$ implies self-duality of
$\Ydg_\cA$. The enhanced symmetry of $\rgrav$ follows from the
skew-symmetry property (Theorem~\ref{V1-thm:thqg-V-collision-twisting}(ii))
applied when $\cA = \cA^!$.
\end{proof}

\begin{remark}[Ordered Yangian structure at the self-dual point]
\label{rem:yangian-vir13-self-dual}
\ClaimStatusProvedHere
\index{Virasoro!self-duality at c=13@self-duality at $c = 13$!ordered Yangian}
\index{gravitational Yangian!self-duality!ordered structure}
%
%: kappa(Vir_c) = c/2, from C2 census; c=13 -> 13/2 verified.
%: self-dual at c=13, NOT c=26.
%
The self-duality at $c = 13$ has a refined expression at the level
of the $\Eone$-ordered bar complex
$\barB^{\Eone}(\mathrm{Vir}_{13})
= T^c(s^{-1}\,\overline{\mathrm{Vir}}_{13})$.
Three structures become simultaneously symmetric.

\emph{(i) Koszul complementarity.}\;
The Koszul conductor for Virasoro satisfies
$K = \kappa(\mathrm{Vir}_c) + \kappa(\mathrm{Vir}_{26-c})
= c/2 + (26-c)/2 = 13$
for all~$c$.
%: kappa(Vir_c) = c/2 from landscape_census.tex; c=0 -> 0, c=13 -> 13/2.
At $c = 13$ the two summands coincide:
$\kappa = \kappa' = 13/2$.
This is the unique central charge at which the
$\Eone$-ordered gravitational Yangian acts on
$\mathrm{Vir}_{13}^{\otimes n}$ with equal weight from both the
algebra and its Koszul dual; away from $c = 13$,
the algebra and dual carry different kappa-weights, and the
$R$-matrix intertwines representations of genuinely different
Virasoro algebras.

\emph{(ii) $R$-matrix involution.}\;
The gravitational $r$-matrix $\rgrav_{\mathrm{Vir}_{13}}(z)
= \sum_{m \geq 0} \Omega_m / z^{m+1}$
acts on $\mathrm{Vir}_{13} \otimes \mathrm{Vir}_{13}$
(not on $\mathrm{Vir}_{13} \otimes \mathrm{Vir}_{26-c}$
with different central charges).
The skew-symmetry
$\rgrav(z)_{12} = -\rgrav(-z)_{21}$
from the preceding proposition becomes an \emph{endomorphism}
property: the $R$-matrix is a genuine involution
$R \colon V \otimes V \to V \otimes V$
on a single representation space $V$, rather than an intertwiner
$V \otimes W \to W \otimes V$ between two different spaces.
This is the ordered ($\Eone$) refinement of the scalar identity
$\kappa = \kappa'$ at degree~$2$.

\emph{(iii) Ordered bar self-duality.}\;
Since $\Ydg_\cA
= \Convstr(\barB^{\Eone}(\cA),\,\chirLie)$
(Remark~\ref{V1-rem:yangian-ordered-boundary-face})
and the Verdier intertwining gives
$\mathbb{D}_{\Ran}(\barB^{\Eone}(\cA))
\simeq \barB^{\Eone}(\cA^!)$,
the self-duality $\cA \cong \cA^!$ at $c = 13$
yields an involutive equivalence
\begin{equation}\label{eq:thqg-V-ordered-bar-self-dual}
\barB^{\Eone}(\mathrm{Vir}_{13})
\;\simeq\;
\mathbb{D}_{\Ran}\bigl(\barB^{\Eone}(\mathrm{Vir}_{13})\bigr)
\end{equation}
of $\Eone$-coalgebras on $\Ran(X)$.
The deconcatenation coproduct is Verdier-self-dual, and the
$R$-matrix data (which the $\Sigma_n$-coinvariant shadow
$\barB^{\Sigma}$ forgets) is preserved by the involution.
Passing to the gravitational Yangian,
$\Ydg_{\mathrm{Vir}_{13}} \cong
\Ydg_{\mathrm{Vir}_{13}^!}$
as dg Lie algebras with $R$-matrix, so the RTT presentation
of the Yangian is invariant under the Koszul involution.

In the gravitational reading: the $c = 13$ Virasoro HT theory
has a graviton-exchange kernel that scatters identical particles
(same central charge on both sides of the cut), the genus-zero
binary shadow is an endomorphism of a single Hilbert space, and
the gravitational Yangian acts by self-intertwining operators.
\end{remark}


% ======================================================================
%
% SUBSECTION 6: LINE-OPERATOR CATEGORIES AND MC3
%
% ======================================================================

\subsection{Line categories and MC3}
% label removed: subsec:thqg-V-line-operators-mc3
\index{line category}
\index{MC3!line categories}

The gravitational Yangian $\Ydg_\cA$ contributes additional structure
to the line side in the HT theory only on the affine/proved
dg-shifted-Yangian lane. Independently, the underlying line category
is modeled by modules for the open-colour dual, written here via the
standard algebraic model
$\cC_\cA^{\mathrm{model}} := \cA^!\text{-}\mathsf{mod}$
\textup{(}Theorem~\textup{\ref{thm:lines_as_modules})}. The MC3
thick-generation programme concerns how much of the additional
Yangian/spectral structure extends from the evaluation-generated core
to richer line-side surfaces. This subsection keeps those two layers
separate.


\subsubsection{The line category and its Yangian enhancement}

\begin{definition}[Line category and conditional Yangian enhancement]
% label removed: def:thqg-V-line-operator-category
\index{line category!from gravitational Yangian}
For a modular Koszul chiral algebra~$\cA$, write
\begin{equation}% label removed: eq:thqg-V-line-category
\cC_\cA^{\mathrm{model}}
\;:=\;
\cA^!\text{-}\mathsf{mod}
\end{equation}
for the module category that models the line category.
On the affine/proved Yangian lane, this line-side model is presented by
dg-shifted-Yangian modules
$\mathrm{Mod}_{\Ydg_\cA}(\mathrm{dgVect})$. When a compatible
dg-shifted-Yangian realization and spectral kernel $r_\cA(z)$ are
supplied, this category carries a meromorphic tensor product
\begin{equation}% label removed: eq:thqg-V-meromorphic-tensor
M \otimes^{\mathrm{mer}}_z N
\;:=\;
\bigl(M \otimes N,\;
d_M \otimes \id + \id \otimes d_N + r_\cA(z)\bigr),
\end{equation}
where $r_\cA(z)$ acts as the $r$-matrix on $M \otimes N$,
producing a family of dg modules parametrized by $z \in \C^*$.

When the same spectral kernel admits a compatible quantization, the
braiding on the model category $\cC_\cA^{\mathrm{model}}$ is the
\emph{$R$-matrix braiding}:
\begin{equation}% label removed: eq:thqg-V-braiding
\beta_{M,N}(z)\colon
M \otimes^{\mathrm{mer}}_z N
\;\xrightarrow{\;\sim\;}
N \otimes^{\mathrm{mer}}_{-z} M,
\qquad
\beta_{M,N}(z) = P_{12} \circ R_\cA(z),
\end{equation}
where $P_{12}$ is the permutation and $R_\cA(z) = 1 + \hbar\,r_\cA(z) +
O(\hbar^2)$ is the quantum $R$-matrix obtained by exponentiating the
classical $r$-matrix. The CYBE
(Theorem~\ref{V1-thm:thqg-V-cybe-from-arnold}) ensures that $\beta$
satisfies the braid relation.
\end{definition}

\begin{remark}[Meromorphic vs.\ genuine braiding]
% label removed: rem:thqg-V-meromorphic-braiding
\index{braiding!meromorphic}
On such a compatible spectral realization, the tensor product on
$\cC_\cA^{\mathrm{model}}$ is \emph{meromorphic} in the spectral
parameter~$z$: it is not defined at $z = 0$ (where the $r$-matrix
has a pole), and the braiding is not a morphism of honest dg modules
but a meromorphic family of morphisms. This is the correct structure
for line operators in HT theories: the spectral parameter is the
transverse holomorphic coordinate, and the singularity at $z = 0$ is
the collision singularity of two line operators at the same point.

The passage from meromorphic braiding to genuine braiding requires
either:
\begin{enumerate}[label=\textup{(\roman*)}]
\item specializing $z = q^h$ for a quantum parameter~$q$ and
 conformal weight~$h$ (the Kazhdan--Lusztig approach,
 Theorem~\ref{V1-thm:kazhdan-lusztig-equivalence}); or
\item taking the monodromy representation of the KZ connection
 around the singularity (the Drinfeld--Kohno approach,
 Theorem~\ref{V1-thm:derived-dk-affine}).
\end{enumerate}
Both approaches produce the same genuine braided tensor category,
which is the category of representations of the quantum group
$U_q(\fg)$ (Drinfeld--Kohno theorem).
\end{remark}


\subsubsection{Thick generation and MC3}

\begin{theorem}[Type-$A$ MC3 reduction via the gravitational Yangian;
\ClaimStatusProvedHere]
% label removed: thm:thqg-V-mc3-thick-generation
\index{MC3!thick generation!gravitational Yangian}
In type~A, the evaluation-generated core on the gravitational
line-side surface is
\begin{equation}% label removed: eq:thqg-V-mc3-generation
\cC^{\mathrm{ev}}_{\widehat{\mathfrak{sl}}_N}
\;:=\;
\mathrm{Thick}\bigl\langle
 V_\lambda(z) \mid \lambda \in P^+(\mathfrak{sl}_N),\;
 z \in \C^*
\bigr\rangle
\;\subset\;
\cC_{\widehat{\mathfrak{sl}}_N}.
\end{equation}
The derived Drinfeld--Kohno comparison is proved on this core by the
Volume~I evaluation-generated-core DK theorem. Beyond the
evaluation-generated core, the Volume~I type-$A$ MC3 reduction theorem
proves the shifted-prefundamental-generation and pro-Weyl-recovery
packets and isolates a single remaining compact-completion gap:
the Volume~I compact-completion conjecture, which asks to extend the
evaluation-core DK equivalence to the compact shifted-prefundamental
core and compare that compact-core equivalence with the completed /
pro-nilpotent line-side category.
\end{theorem}

\begin{proof}
The definition of $\cC^{\mathrm{ev}}_{\widehat{\mathfrak{sl}}_N}$ is
tautological. The Volume~I evaluation-generated-core DK theorem proves
the DK comparison on this core. The Volume~I type-$A$ MC3 reduction theorem
then proves the shifted-prefundamental-generation and pro-Weyl
packages and states that the only remaining post-core step is
the compact-completion conjecture. Transporting those
results to the line-side surface gives the claim.
\end{proof}

\begin{corollary}[Type-$A$ DK-5 reduction to the compact-completion packet;
\ClaimStatusProvedHere]
% label removed: cor:thqg-V-dk5-type-a
\index{Drinfeld--Kohno!DK-5!type A accessibility}
On the gravitational line-side surface for
$\widehat{\mathfrak{sl}}_N$, the only remaining obstruction between the
proved evaluation-core DK comparison and a full DK-5 statement is the
compact-completion packet of
the Volume~I compact-completion conjecture. No further independent
type-$A$ obstruction remains beyond that packet.
\end{corollary}

\begin{proof}
Immediate from the preceding theorem and the Volume~I type-$A$ MC3
reduction theorem.
\end{proof}

\begin{remark}[DK-5 resolved for $\mathfrak{sl}_2$]
\label{rem:vol2-dk5-sl2}
\index{Drinfeld--Kohno!DK-5!sl2@$\mathfrak{sl}_2$ resolved}
For $\fg = \mathfrak{sl}_2$, the full chain from the bar complex to the
Verlinde category is now complete
\textup{(}Volume~I,
Proposition~\textup{\ref*{V1-prop:dk5-sl2-frt})}: the ordered bar
complex $\barB^{\mathrm{ord}}(\widehat{\mathfrak{sl}}_{2,k})$ produces
the Yang $R$-matrix $R(u) = uI + \hbar P$, the FRT construction
recovers $U_q(\mathfrak{sl}_2)$ with its Hopf algebra structure, and
at roots of unity $q = e^{2\pi i/(k+2)}$ the Verlinde truncation follows
from the vanishing quantum dimension $[k{+}2]_q = 0$. This is the first
instance where the entire DK-5 programme is verified end-to-end.
\end{remark}

\begin{remark}[MC3 for all simple types: evaluation core proved, post-core conditional]
% label removed: rem:thqg-V-mc3-arbitrary-type
\index{MC3!all simple types}
The prefundamental Clebsch--Gordan closure
\textup{(}package~\textup{(i)} of MC3\textup{)} is proved for all simple
types by the Volume~I all-types categorical Clebsch--Gordan theorem,
and the DK comparison on the evaluation-generated core is proved for
all simple types by the corresponding Volume~I corollary. In type~$A$,
the Volume~I type-$A$ MC3 reduction theorem proves the
shifted-prefundamental-generation and pro-Weyl-recovery packets,
leaving only the compact-completion extension conjecture. For other
types, the analogous post-core mechanism becomes type-independent only
conditionally, once the shifted-prefundamental, pro-Weyl, and
compact-completion packets of
the conditional Volume~I automatic-generalization mechanism are separately
supplied in that type.
\end{remark}

\begin{construction}[MC3 landscape: all simple types]
% label removed: constr:vol2-mc3-gradient
The all-types theorematic surface consists of categorical
Clebsch--Gordan closure and the evaluation-generated-core DK
comparison from Volume~I. In type~$A$, the Volume~I type-$A$ MC3 reduction theorem
proves the shifted-prefundamental-generation and pro-Weyl-recovery
packets, leaving only the compact-completion extension conjecture.
For arbitrary type, the conditional automatic-generalization mechanism
shows that the same post-core mechanism is available only
conditionally, once the corresponding shifted-prefundamental,
pro-Weyl, and compact-completion packets are supplied in that type.
\end{construction}

\begin{remark}[Type $B_2$ after categorical CG closure]%
% label removed: conj:mc3-type-B2
\index{MC3!type B2@type $B_2$}
For $\fg = \mathfrak{so}_5$ \textup{(}type $B_2$\textup{)}, the
categorical prefundamental Clebsch--Gordan decomposition
\begin{equation}% label removed: eq:mc3-B2-cg
V_{\omega_i}(a) \otimes L^-_i(b)
\;\cong\;
\bigoplus_j L^-_i(b_j)
\end{equation}
holds as $Y(\mathfrak{so}_5)$-modules in
$\cO^{\mathrm{sh}}_{\leq 0}$, not merely at the $K_0$-level.
This is a special case of the all-types result
\textup{(}proved in Volume~I on the evaluation-generated core\textup{)}.

Promoting this evaluation-core statement to the full post-core MC3
package for $\mathfrak{so}_5$ would require the
shifted-prefundamental, pro-Weyl, and compact-completion packets of
the conditional Volume~I automatic-generalization mechanism. Those packets
are not proved here for type~$B_2$.

\smallskip\noindent
\textbf{Proof method.} The uniform all-types proof
in Volume~I uses the
multiplicity-free $\ell$-weight property (Chari--Moura, 2006)
to bypass the folding strategy entirely. Type $B_2$ has two
fundamental weights $\omega_1$ \textup{(}vector\textup{)} and
$\omega_2$ \textup{(}spinor\textup{)}. The CG closure at
$K_0$-level
\textup{(}the Volume~I prefundamental Clebsch--Gordan proposition\textup{)}
lifts categorically via multiplicity-freeness of $\ell$-weight spaces.
\end{remark}


\subsubsection{The Yangian-shadow principle}

\begin{principle}[Yangian-shadow principle]
% label removed: princ:thqg-V-yangian-shadow
\index{Yangian-shadow principle}
For every holomorphic-topological theory~$T$ with boundary chiral
algebra~$\cA$:
\begin{enumerate}[label=\textup{(\Alph*)}]
\item The genus-zero binary shadow
 $r_\cA(z) = \Rescoll_{0,2}(\Theta_\cA)$ recovers the
 dg-shifted-Yangian data on the affine / evaluation-core lane;
 independently, a standard algebraic model for the line side is
 $\cC_\cA^{\mathrm{model}} := \cA^!\text{-}\mathsf{mod}$.
\item The genus-zero $n$-point shadow
 $\Sh_{0,n}(\Theta_\cA)$ identifies on the affine Kac--Moody
 comparison surface with the KZ connection, and on the Virasoro
 degenerate-representation comparison surface its higher collision
 residues yield the BPZ null-vector equations.
\item The formal genus-refined binary shadow series
 $R^{\mathrm{bin}}_\cA(z;\hbar) =
 \sum_g \hbar^{2g}\,r_{\cA,g}(z)$ is the candidate higher-genus
 binary package extracted from the shadow coefficients.
\item On the conjectural modular-Yangian surface, a genuine modular
 kernel in $\Ymod_\cA$ would satisfy the stable-graph YBE and furnish
 the higher-genus completion of the genus-zero Yangian package.
\end{enumerate}
Components~(A)--(B) are proved
\textup{(}Theorems~\textup{\ref{V1-thm:collision-residue-twisting}},
\textup{\ref{V1-thm:collision-depth-2-ybe}},
and the Volume~I affine shadow/KZ comparison theorem\textup{)}.
Component~(C) is the proved formal binary-shadow series of
Theorem~\textup{\ref*{V1-thm:thqg-V-pro-mc-element}}.
Component~(D) is frontier modular-kernel data, not a theorem of this
chapter.

The principle asserts: \emph{the gravitational Yangian is not a new
structure imposed on the HT theory; its universal proved core is the
genus-zero binary shadow of the universal MC element, while the
higher-genus modular completion is an intended extension of that
shadow package.}
\end{principle}


\subsubsection{Dictionary: Yangian data vs.\ gravitational data}

\begin{remark}[Yangian-gravity dictionary]
% label removed: rem:thqg-V-dictionary
\index{gravitational Yangian!dictionary}
The following table translates between the Yangian language of
\S\ref{V1-chap:yangians} and the gravitational language of the present
chapter:

\begin{center}
\renewcommand{\arraystretch}{1.3}
\small
\begin{tabular}{@{}lll@{}}
\toprule
\textbf{Yangian object} & \textbf{Gravitational object}
 & \textbf{Status} \\
\midrule
$r_\cA(z)$ & Binary graviton exchange in the gravitational reading
 & Proved (Thm~\ref{V1-thm:thqg-V-collision-twisting}) \\
CYBE & Tree-level $3$-graviton consistency in the gravitational reading
 & Proved (Thm~\ref{V1-thm:thqg-V-cybe-from-arnold}) \\
$\Ydg_\cA$ & Genus-zero scattering algebra on the dg-shifted-Yangian lane
 & Proved (Def~\ref{V1-def:thqg-V-modular-yangian}) \\
$R^{\mathrm{bin}}_\cA(z;\hbar)$ & Formal higher-genus binary shadow series
 & Proved (Thm~\ref*{V1-thm:thqg-V-pro-mc-element}) \\
sgYBE & All-loop graviton consistency in the modular-kernel reading
 & Conjectural modular-kernel surface
 (Def~\ref*{V1-def:thqg-V-stable-graph-ybe}) \\
$\cC_\cA^{\mathrm{model}} := \cA^!\text{-mod}$ & Module model for the line category
 & Proved (Thm~\ref{thm:lines_as_modules}) \\
RTT relation & Evaluation-core graviton factorization in the gravitational reading
 & Proved (Prop~\ref{V1-prop:thqg-V-rtt-from-sgybe}) \\
Evaluation module $V_\lambda(z)$ & Evaluation-core line object
 & Proved on the evaluation-generated core
 (Volume~I) \\
Shadow connection $\nabla^{\mathrm{hol}}$
 & KZ on affine Kac--Moody comparison surface; BPZ on Virasoro
 degenerate-representation surface
 & Proved (Volume~I affine shadow/KZ comparison theorem and
 Virasoro/BPZ comparison proposition) \\
Quartic contact $\mathfrak{Q}$ & $4$-graviton obstruction in the gravitational reading
 & Proved (Comp~\ref{V1-comp:thqg-V-quartic-graviton}) \\
\bottomrule
\end{tabular}
\end{center}
\emph{The proved entries are direct projections of the single MC
element $\Theta_\cA$; the sgYBE row records the conjectural
higher-genus completion surface.}
\end{remark}


\subsubsection{The gravitational Yangian as ordered boundary face}

\begin{remark}[The gravitational Yangian as $\Eone$-ordered face;
cf.~Remark~\textup{\ref{V1-rem:yangian-ordered-boundary-face}}]
% label removed: rem:thqg-V-ordered-face
\index{gravitational Yangian!as ordered boundary face}
The gravitational Yangian $\Ydg_\cA$ is the $\Eone$-ordered face of
the universal twisting:
\begin{equation}% label removed: eq:thqg-V-ordered-face
\Ydg_\cA
\;=\;
\Convstr\!\bigl(\barB^{\Eone}(\cA),\;\chirLie\bigr),
\end{equation}
the strict convolution dg Lie algebra of the $\Eone$-bar complex.
The bar complex $\barB^{\Eone}(\cA)$ lives on \emph{ordered}
configurations, the Verdier duality inverts the $R$-matrix
($\mathbb{D}\,r(z) = r(-z)^{-1}$, not the level), and the shadow
tower on ordered configurations encodes the full spectral-parameter
dependence of the rational $R$-matrix through all degrees.

The passage from the unordered ($\Einf$-chiral) modular
convolution algebra $\gAmod$ to the ordered ($\Eone$-chiral)
gravitational Yangian $\Ydg_\cA$ is not a restriction but a
refinement: the ordered structure retains the spectral-parameter
dependence that the unordered structure averages over. For
commutative ($\Einf$-chiral) algebras such as $\widehat{\fg}_k$,
the gravitational Yangian is the Yang $R$-matrix; for genuinely
$\Eone$-chiral algebras such as $Y(\fg)$, it is the full RTT
Yangian.
\end{remark}


\subsubsection{Examples: gravitational Yangians of the standard families}

\begin{computation}[Gravitational Yangian of Heisenberg;
\ClaimStatusProvedHere]
% label removed: comp:thqg-V-heis-yangian
\index{Heisenberg algebra!gravitational Yangian}
For rank-$N$ Heisenberg $\mathcal{H}^N_\kappa$:
\begin{enumerate}[label=\textup{(\roman*)}]
\item $\Ydg_{\mathcal{H}} = \C[\![z^{-1}]\!]$
 (abelian, rank $1$).
\item $R^{\mathrm{bin}}_{\mathcal{H}}(z;\hbar) = N/(\kappa z)$
 (no higher-genus binary corrections).
\item A standard model for the line side is
 $\cC_{\mathcal{H}}^{\mathrm{model}} = \mathrm{Vect}^{\Z}$
 (graded vector spaces, one line object per charge).
\item The higher-genus binary shadow corrections vanish, so the
 conjectural modular completion would add no new binary terms.
\end{enumerate}
The Heisenberg gravitational Yangian is the simplest possible
case: one generator, one relation (which is trivial), no quantum
corrections. In the gravitational reading, this corresponds to a free
theory.
\end{computation}

\begin{computation}[Gravitational Yangian of affine KM;
\ClaimStatusProvedHere]
% label removed: comp:thqg-V-affine-yangian
\index{Kac--Moody!gravitational Yangian}
For $\widehat{\fg}_k$ at non-critical level:
\begin{enumerate}[label=\textup{(\roman*)}]
\item $\Ydg_{\widehat{\fg}_k} \cong Y(\fg)$
 (the standard Yangian of~$\fg$).
\item $r_{\widehat{\fg}_k,0}(z) = \Omega_\fg/((k{+}h^\vee)\,z)$
 (Casimir $r$-matrix).
\item A standard model for the line side is
 $\cC_{\widehat{\fg}_k}^{\mathrm{model}} \simeq
 \cO^{\mathrm{sh}}_q(\widehat{\fg})$
 at $q = e^{i\pi/(k+h^\vee)}$.
\item At genus~$0$, the collision/Yang--Baxter package is the
 standard CYBE for $\fg$; the corresponding higher-genus modular
 completion is expected to involve KZB-type corrections.
\end{enumerate}
The affine KM gravitational Yangian is the classical Yangian~$Y(\fg)$:
in the gravitational reading, this corresponds to the semistrict
higher-spin theory whose scattering is read through the rational
$R$-matrix of~$\fg$.
\end{computation}

\begin{computation}[Gravitational Yangian of $\mathrm{Vir}_c$;
\ClaimStatusProvedHere]
% label removed: comp:thqg-V-vir-yangian-summary
\index{Virasoro!gravitational Yangian!summary}
For $\mathrm{Vir}_c$:
\begin{enumerate}[label=\textup{(\roman*)}]
\item $\Ydg_{\mathrm{Vir}_c}$ is an infinite-dimensional
 dg Lie algebra with generators $\{\Omega_m\}_{m \geq 0}$
 (Construction~\ref{constr:thqg-V-vir-r-explicit}).
\item $\rgrav_{\mathrm{Vir}}(z) =
 \sum_{m \geq 0} \Omega_m/z^{m+1}$ (infinite series of
 Casimir components).
\item A standard model for the line side is
 $\cC_{\mathrm{Vir}_c}^{\mathrm{model}} =
 \mathrm{Vir}_{26-c}\text{-}\mathsf{mod}$
 (the expected Virasoro module model for the line side).
\item The candidate higher-genus binary-shadow package suggests an
 infinite tower of corrections beyond the genus-zero CYBE.
\item Self-duality at $c = 13$; degeneration at $c = 0$;
 the depth-zero resonance shadow at $c = 26$.
\item The quartic contact invariant
 $Q^{\mathrm{contact}}_{\mathrm{Vir}} = 10/[c(5c+22)]$
 is the first obstruction to the $r$-matrix being a
 complete description.
\end{enumerate}
The Virasoro gravitational Yangian encodes the richest explicit
genus-zero binary shadow data for HT theories with stress-tensor
boundary operators. In the gravitational reading, its infinite
shadow obstruction tower is expected to model an infinite graviton-vertex tower,
while the full higher-genus modular completion remains frontier.
\end{computation}


\subsubsection{The gravitational completion problem}

\begin{remark}[Gravitational completion and MC4]
% label removed: rem:thqg-V-gravitational-completion
\index{gravitational Yangian!completion problem}
\index{MC4!gravitational completion}
The MC4 completion theorem
(Theorem~\ref{V1-thm:completed-bar-cobar-strong}) ensures that the
gravitational Yangian admits a completion:
\begin{equation}% label removed: eq:thqg-V-completion
\widehat{\Ydg}_\cA
\;:=\;
\varprojlim_{N}\;
\Ydg_\cA / F^{N+1}\Ydg_\cA,
\end{equation}
with the strong filtration axiom
$\mu_r(F^{i_1},\dotsc,F^{i_r}) \subseteq F^{i_1+\cdots+i_r}$
ensuring that the completion is compatible with the $A_\infty$
structure (Lemma~\ref{V1-lem:degree-cutoff}).

For the Virasoro gravitational Yangian: the completion
$\widehat{\Ydg}_{\mathrm{Vir}}$ is the pronilpotent envelope of
the infinite-dimensional dg Lie algebra generated by the Casimir
components $\{\Omega_m\}_{m \geq 0}$. The MC4 splitting
distinguishes:
\begin{enumerate}[label=\textup{(\roman*)}]
\item \emph{MC4$^+$ (positive towers)}:
 $\mathcal{W}_{1+\infty}$, affine Yangians, solved by weight
 stabilization
 (Theorem~\ref{V1-thm:stabilized-completion-positive}).
\item \emph{MC4$^0$ (resonant towers)}:
 Virasoro, non-quadratic $\mathcal{W}_N$, reduced to finite
 resonance by
 Theorem~\ref{V1-thm:resonance-filtered-bar-cobar}.
\end{enumerate}
The gravitational completion problem is: identify the completed
gravitational Yangian $\widehat{\Ydg}_\cA$ with a standard
quantum group. For affine KM, this is the quantum group
$U_q(\fg)$ (Drinfeld--Kohno). For Virasoro, the target is the
$\mathcal{W}$-algebra quantum group, whose identification is
part of the MC4 resonance programme.
\end{remark}


\subsubsection{Summary: the gravitational Yangian as projection of $\Theta_\cA$}

\begin{remark}[Result G5: synthesis]
% label removed: rem:thqg-V-g5-synthesis
\index{gravitational Yangian!synthesis}
Result~G5 is the statement that the gravitational Yangian
$\Ydg_\cA$ and its ambient pronilpotent modular completion
$\Ymod_\cA$ are organized by the single MC element $\Theta_\cA$
(Theorem~\ref*{V1-thm:mc2-bar-intrinsic}). The content is:
\begin{enumerate}[label=\textup{(\arabic*)}]
\item The collision filtration on $\gAmod$ stratifies the
 genus-zero sector by FM boundary depth.
\item The binary collision residue
 $\Rescoll_{0,2}(\Theta_\cA) = \pi_\cA$ recovers the
 universal twisting morphism and the spectral $r$-matrix.
\item The Arnold relation on $\overline{\mathcal{M}}_{0,4}$ forces
 the CYBE\@.
\item The pronilpotent completion $\Ymod_\cA$ provides the ambient
 modular dg Lie package, and the genus-refined binary projections
 define the formal series
 $R^{\mathrm{bin}}_\cA(z;\hbar)=\sum_g \hbar^{2g}r_{\cA,g}(z)$.
 Comparison with a genuine modular kernel
 $R^{\mathrm{mod}}_\cA(z;\hbar)$ satisfying the sgYBE is frontier.
\item For Virasoro, the binary kernel $\rgrav_{\mathrm{Vir}}(z)$ has higher-order
 poles encoding the infinite shadow obstruction tower; the quartic contact
 invariant $Q^{\mathrm{contact}} = 10/[c(5c+22)]$ is the first
 nonlinear obstruction; and self-duality holds at $c = 13$.
\item The line category is modeled here by
 $\cC_\cA^{\mathrm{model}} := \cA^!\text{-}\mathsf{mod}$,
 and this line-side model is proved independently,
 while the Yangian MC3 surface has an all-types evaluation-core
 theorem and a post-core compact/completion frontier.
\end{enumerate}
Items~(1)--(3), together with the ambient completion in~(4) and the
line-category statement in~(6), are consequences of the proved
algebraic core. The gravitational Yangian is not an additional
structure; its universal proved content is the ordered binary shadow
of the universal MC class, while the modular-kernel completion in~(4)
remains frontier.
\end{remark}

\subsection{Gravitational Yangian as
$Q_\fg^{k,f,\mu}|_{\fg=\mathrm{Vir}}$}
\label{subsec:thqg-gravitational-yangian-anchor}

\begin{remark}[Gravitational case as type-2 dg-shifted on the
infinite-dimensional Virasoro algebra; \ClaimStatusProvedHere]
\label{rem:thqg-gravitational-yangian-yangian-type}
In the four-Yangian-types disambiguation, the gravitational
dg-shifted Yangian $Y^{\mathrm{dg}}(\mathrm{Vir}_c)$ is the
\emph{type-2} object (complete filtered $A_\infty$-algebra at a
point/formal disk) with
$\fg$ replaced by the infinite-dimensional Virasoro algebra
$\mathrm{Vir}_c$. The collapse phenomenon already noted in this
chapter (the level-$r$ generators $\mathsf{T}_n^{(r)} = \binom{n+r-1}{r}\,L_{-n-r}$ are scalar multiples of single
Virasoro modes; the underlying Lie reduces to
$\mathrm{Vir}_{26-c}$) does not eliminate the dg content: the
$A_\infty$ tower $\{m_k\}_{k\ge 3}$ remains nontrivial for generic
$c$ and is the genuine output of the construction. As a
specialisation of the unified
$Q_\fg^{k,f,\mu}$ (Theorem~\ref{thm:unified-chiral-quantum-group}),
the gravitational Yangian is the fibre at $(\fg, k, f, \mu) =
(\mathrm{sl}_2, k_{\mathrm{Sug}}, f_{\mathrm{prin}}, 0)$ with
$c = c(k_{\mathrm{Sug}}) = 1 - 6(k+2)^{-1}(k+1)^2$ via the
Virasoro Sugawara construction.
\end{remark}

\begin{remark}[Notion-B for the gravitational dg-shifted Yangian;
\ClaimStatusProvedElsewhere]
\label{rem:thqg-gravitational-yangian-notion-b}
$Y^{\mathrm{dg}}(\mathrm{Vir}_c)$ is read as a Notion-B
$E_1$-chiral algebra: structure maps $\{m_k\}_{k\ge 1}$ in
$\mathrm{End}^{\mathrm{ch}}_{\mathrm{Vir}_c}$, Stasheff identities
chain-level, the binary residue $r^{\mathrm{grav}}_1 = (c/2)
L_{-1}^{\otimes 2}$ as the leading-order Maurer--Cartan datum.
The $\Sigma_2$-coinvariant projection
$\mathrm{av}(r^{\mathrm{grav}}(z)) = \kappa(\mathrm{Vir}_c) = c/2$
recovers the modular characteristic at degree~$2$, in agreement
with the $E_1$-primacy directive.
\end{remark}

\begin{remark}[Level-prefix discipline for the gravitational kernel;
\ClaimStatusProvedElsewhere]
\label{rem:thqg-gravitational-yangian-rmatrix}
The gravitational $r$-matrix $r^{\mathrm{grav}}(z) = \sum_m \Omega_m
z^{-m-1}$ carries the central charge $c$ as its level prefix
implicitly through the binary residue normalization
$\Omega_0 = (c/2)\,L_{-1}^{\otimes 2}$. At $c=0$ the
gravitational scattering vanishes (the formal-series identity must
be specified algebraically); at $c=13$ the Koszul-self-dual point
gives the gravitational Yangian its self-dual specialisation; at
$c=26$ the matter--ghost critical locus emerges. The level prefix
$c$ is load-bearing and never omitted in this chapter (verified
verbatim against the gravitational kernel formulae above).
\end{remark}

\begin{remark}[Miura $(\Psi-1)/\Psi$ universality on the spectral
coproduct of the gravitational Yangian; \ClaimStatusProvedElsewhere]
\label{rem:thqg-gravitational-yangian-miura-coproduct}
The spectral coproduct
$\Delta_z\colon \mathrm{Vir}_c \to (\mathrm{Vir}_c)^{\otimes 2}((z))$
of the gravitational dg-shifted Yangian has the form
\[
\Delta_z(T) \;=\; T\otimes 1 \;+\; 1\otimes T \;+\;
\frac{\Psi-1}{\Psi}\,(J\otimes J)\,z^{-1} \;+\; O(z^{-2})
\]
with $\Psi = c/12$ being the Drinfeld--Sokolov coupling reading.
The $(\Psi-1)/\Psi$ coefficient is the universal Miura
cross-term universality (Vol~I
Theorem~\ref*{V1-thm:miura-cross-universality}). At
$\Psi=2$ ($c=24$, the Schellekens point) the substitution gives
$1/\Psi = (\Psi-1)/\Psi = 1/2$, an algebraic coincidence; the
correct universality is $(\Psi-1)/\Psi$ for every spin $s\ge 2$,
verified at $\Psi=3$ ($c=36$).
\end{remark}

\begin{remark}[Gravitational Yangian as boundary input to the
universal-holography master theorem]
\label{rem:thqg-gravitational-yangian-uhm}
The gravitational dg-shifted Yangian
$Y^{\mathrm{dg}}(\mathrm{Vir}_c)$ is the open-colour boundary
input to \cref{thm:universal-holography-master} (Vol~II
\cref{ch:programme-climax-platonic}) at the class M (Virasoro)
specialisation of $\Phi_{\mathrm{hol}}$: the bulk 3d HT theory
$T_{\mathrm{Vir}_c}$ is the chiral derived centre
$\mathrm{C}^\bullet_{\mathrm{ch}}(\mathrm{Vir}_c, \mathrm{Vir}_c)$,
identified via the Drinfeld--Sokolov/Hochschild compatibility bridge
(\cref{thm:chd-ds-hochschild}) with
$H^\bullet_{\mathrm{DS}}(\mathrm{C}^\bullet_{\mathrm{ch}}
(V_k(\mathrm{sl}_2), V_k(\mathrm{sl}_2)))$ at the level $k$ giving
the desired $c$ via Sugawara. The Vol~I unified Koszul reflection
is the bar-side input on the boundary; the Vol~I universal-conductor
theorem controls the dual reflections $c\mapsto 13-c$ (Koszul) and
$c\mapsto 26-c$ (matter--ghost), which agree at the joint fixed
point $c=13$ (self-dual Koszul Virasoro) only via the
gravitational-Yangian collapse seen in this chapter.
\end{remark>
